\documentclass[a4paper,11pt, titlepage, twoside]{article}

%Cedit a tots els estudiants futurs
\title{Estudi i desenvolupament d'una
arquitectura d'intel·ligència artificial
neuromòrfica basada en xarxes
neuronals de polsos (SNNs).}
\author{Alejandro Pérez Hernández}
\date{gener 2026}


% Language and font encodings
\usepackage[catalan]{babel} %Idioma del document
\usepackage{lipsum}
\usepackage[utf8]{inputenc}
\usepackage{lmodern}
\usepackage[T1]{fontenc}
\usepackage[autostyle=true, style=spanish]{csquotes}
\usepackage{parskip}
\setlength{\parskip}{4mm}
\setlength{\footskip}{60pt}
\setlength{\headheight}{15pt}

%% Sets page size and margins
\usepackage[a4paper,top=3cm, bottom=3.2cm, inner=3cm, outer=2cm, footnotesep=1cm, heightrounded]{geometry}


%% Useful packages

\usepackage{listings}
\usepackage{xcolor}
\lstdefinestyle{logs}{
    basicstyle=\ttfamily\small,
    breaklines=true,
    frame=single,
    columns=fullflexible,
    keepspaces=true,
    showstringspaces=false
	captionpos=b
}
% Colors
\definecolor{pyfunction}{HTML}{0000CC} % funcions Python
\definecolor{torchcolor}{HTML}{C45A00}  % torch / PyTorch
\definecolor{varcolor}{HTML}{005F99}    % variables
\definecolor{classfunction}{HTML}{B8860B} % groc fosc llegible per a classes i funcions definides per l'usuari

\lstdefinestyle{codi}{
    basicstyle=\ttfamily\small,
    breaklines=true,
    frame=single,
    columns=fullflexible,
    keepspaces=true,
    showstringspaces=false,
    captionpos=b,
    language=Python,
    numbers=left,
    numberstyle=\tiny,
    stepnumber=1,
    numbersep=8pt,
	commentstyle=\color{green!50!black}\itshape,
    emph={class,def,print,for,if,range,len,open,input,int,float,str,bool,list,dict,set,tuple,zip,enumerate,map,filter,sum,min,max,abs,round,sorted,type,isinstance},
    emphstyle=\color{pyfunction},
    emph={[2]torch,rand_like,Tensor,tensor,randn,rand,zeros,ones,empty,arange,linspace,cat,stack,reshape,view,permute,transpose,matmul,mm,bmm,mean,argmax,argmin,softmax,sigmoid,relu,tanh,no_grad,inference_mode,save,load,nn,Module,Linear,Conv2d,ReLU,Sequential,CrossEntropyLoss,MSELoss,optim,Adam,SGD,AdamW,DataLoader,Dataset},
    emphstyle={[2]\color{torchcolor}\bfseries},
    emph={[3]x,y,X,Y,model,optimizer,criterion,loss,epoch,batch,output,target,logits,pred,accuracy,device,lr},
    emphstyle={[3]\color{varcolor}},
	emph={[4]__init__,__call__,__str__,__repr__,__len__,__getitem__,__setitem__,__iter__,__next__,__enter__,__exit__,forward, in},
	emphstyle={[4]\color{classfunction}\bfseries}
}
\usepackage{float}
\usepackage{verbatim}
\usepackage{amsmath}
\usepackage{systeme}
\usepackage{graphicx}
\graphicspath{{figures/}}
\usepackage{multirow}
\usepackage{caption}
\captionsetup[lstlisting]{position=bottom}
\usepackage{subcaption}
\usepackage{wrapfig}
\usepackage[colorinlistoftodos]{todonotes}
\usepackage{hyperref}
\hypersetup{
    colorlinks=true,
    linkcolor=black, % Enllaços interns: seccions, figures, taules, ToC...
    citecolor=blue,  % Citacions bibliogràfiques
    urlcolor=blue    % URLs externes
}
\usepackage{titlesec}
\usepackage{fancyhdr}
\usepackage[firstpage]{draftwatermark}
\usepackage{transparent}
\usepackage{textcomp} %podem escriure ``o'' amb la comanda \textdegree també € amb \texteuro
\usepackage[gen]{eurosym} %tb official en lloc de ``gen''
\usepackage[section]{placeins} %Evita que les figures saltin de secció
\usepackage{fancyref}
\usepackage{array}
\usepackage{longtable}
\usepackage{mathtools}
\usepackage{commath}
\usepackage{scrextend}%Indentar un bloc sencer
\usepackage{tgpagella}
\usepackage{dtk-logos}
\usepackage{appendix}
\usepackage[backend=biber,style=ieee]{biblatex}
\setlength{\bibitemsep}{0.6em}

\addbibresource{Bibliografia.bib} % nom del teu fitxer .bib

%%Comentaris
\begin{comment}
Aquí podem posar tots els comentaris que calgui sense anar posant %
REFERENCIAT
\label{marker}, \ref{marker} and \pageref{marker} \footnote{footnote text}

\begin{addmargin}[esq]{dta}
El text d'aquí incrementa els seus marges segons ``esq'' i ``dta''
\end{addmargin}
\end{comment}

%%Comença pàgina nova per cada Secció SI hi ha alguna cosa escrita a la anterior

\newcommand{\sectionbreak}{\cleardoublepage}


%%Caps de pagina i peus (E/O (even/odd), L/C/R (left/center/right) y H/F (header/footer))
\fancyhf{}
\fancyhead[ER]{Memòria}
\fancyhead[OL]{Estudi i desenvolupament d'una arquitectura d'intel·ligència artificial neuromòrfica basada en xarxes neuronals de polsos (SNNs).}
\fancyhf[ELH, ORH]{pàg. \thepage}
%\fancyfoot[C]{}
\fancyfoot[EL, OR]{

\includegraphics[scale=0.2]{ETSEIB.png}}
\pagestyle{fancy}
\raggedbottom

\SetWatermarkText{\hspace{9mm}\transparent{0.15}
\includegraphics[scale=2.5]{ETSEIB.png}}
\SetWatermarkAngle{0}
\SetWatermarkLightness{1}



\begin{document}
\renewcommand{\refname}{Bibliografia}
\begin{titlepage}
	{\centering
		{\Huge Treball de Fi de Grau}\\
		\vspace{5mm}
		{\Large \textbf{Grau en Enginyeria en Tecnologies Industrials (GETI)}}\\
		\vspace{20mm}
		\Huge \textbf{Estudi i desenvolupament d'una arquitectura d'intel·ligència artificial neuromòrfica basada en xarxes neuronals de polsos (SNNs).}\\
		\vspace{10mm}
		\Huge\textbf{MEMÒRIA}\\
		\vspace{3mm}
		\Large\today\\  %Si e lloc del dia de la darrera edició es vol una data fixa, elimineu \today i poseu la data
	}
	\vspace{20mm}
	\hspace{2mm}
	\begin{tabular}{l@{ } l}
		\vspace{5mm}
		\Large \textbf{Autor:}        & \Large{Alejandro Pérez Hernández} \\
		\vspace{5mm}
		\Large\textbf{Director:}      & \Large{Álvaro Gómez Pau}          \\

		\Large\textbf{Convocatòria: } & \Large{Abril 2026}                \\
	\end{tabular}\par
	\vspace{10mm}
	{\centering
	
\includegraphics[scale=0.3]{ETSEIB.png}\\
	{\Large Escola Tècnica Superior \\ d'Enginyeria Industrial de Barcelona}\\
	\vspace{3mm}
	
\includegraphics[scale=0.4]{UPC_logo.PNG}
	\par
	}
\end{titlepage}
%\maketitle

\clearpage
\thispagestyle{empty}
\null\newpage
\pagenumbering{arabic}

\section*{Resum}
Aquest treball estudia el funcionament i el potencial de les xarxes neuronals de polsos, o \textit{Spiking Neural Networks} (SNNs), com a arquitectura d’intel·ligència artificial neuromòrfica inspirada en el comportament del cervell humà.

El treball comença presentant l'estat de l'art de la intel·ligència artificial, amb un marc històric i tècnic que presenta des dels primers models de neurona artificial fins a les arquitectures modernes de deep learning. Es descriu l’evolució de les xarxes neuronals artificials, les seves principals aplicacions i les limitacions que presenten en l’actualitat, especialment pel que fa al consum computacional i energètic dels models de gran escala. En base a això, s’introdueixen les SNNs com una possible evolució de les ANNs, ja que incorporen de manera explícita la dimensió temporal i utilitzen un processament basat en esdeveniments, més proper al funcionament biològic de les neurones, el que podria otorgar-les una ventaja en termes de eficiència energètica i capacitat de processament temporal.

A continuació, s’analitzen els fonaments tècnics i matemàtics de les xarxes neuronals de polsos. Es descriuen models neuronals existents, com el de Hodgkin-Huxley, Integrate-and-Fire i Leaky Integrate-and-Fire, així com el paper de les sinapsis, els retards temporals i les diferents estratègies de codificació de la informació. També s’examinen els principals mètodes d’entrenament de les SNNs, incloent-hi l’aprenentatge local basat en STDP, l’ús de gradients substitutius i la conversió d’una ANN prèviament entrenada a una SNN. Aquesta última estratègia serà la utilitzada en la part pràctica del treball, doncs permet aprofitar la maduresa i relativa facilitat dels mètodes d’entrenament de les ANNs i traslladar els pesos apresos a un model de polsos.

La part experimental, exposada a continuació, es divideix en dos casos d’estudi. En el primer, es desenvolupa una ANN senzilla per classificar punts, segons si es troben dins o fora d’un polígon generat de manera pseudoaleatòria. Un cop entrenada, aquesta xarxa es converteix en una SNN mitjançant codificació per taxa (\textit{Rate Encoding}), normalització de pesos i simulació temporal de polsos. Els resultats mostren que la ANN assoleix una precisió del 93,0\%, mentre que la SNN convertida arriba al 90,1\%, amb una pèrdua de rendiment relativament petita. Aquests resultats confirmen que la conversió ANN-SNN es viable i pot conservar de manera raonablement satisfactòria el comportament del model original d'ANN.

En el segon cas d’estudi, s’aplica el mateix plantejament al dataset MNIST, un dataset format per imatges de dígits manuscrits. En aquest cas, es parteix d’una xarxa neuronal convolucional (CNN) entrenada per classificar imatges del 0 al 9, que posteriorment es converteix en una SNN. La CNN obté una exactitud del 99,06\%, mentre que la SNN convertida arriba al 93,19\%. Tot i que la pèrdua de precisió és més elevada que en el cas geomètric, els resultats demostren que la conversió també és viable en un problema de classificació d’imatges, una casuística que s'assimila més a una aplicació real. Com a complement per visualitzar els resultats obtinguts, s’ha desenvolupat una aplicació web que permet dibuixar dígits manualment i comparar les prediccions de la CNN i la SNN.

Finalment, el treball analitza el possible impacte econòmic, ambiental i social que l'implementació d'aquests tipus de xarxes podrien tenir. Les SNNs no es plantegen com una substitució immediata de les xarxes neuronals convencionals, sinó com una línia tecnològica amb un gran potencial en escenaris on el consum energètic és crític. La seva naturalesa basada en esdeveniments pot reduir l’activitat computacional i, per tant, el consum elèctric i la calor generada, fet que podria contribuir a disminuir els costos operatius i l’impacte ambiental dels sistemes d’intel·ligència artificial. Els resultats obtinguts indiquen doncs que les SNNs són una alternativa prometedora, especialment si es combinen amb processos de refinament específics i amb hardware neuromòrfic capaç d’explotar les seves avantatges.

%----------------------------------------------------------------------------------------------
\section*{Abstract}
This project studies the functioning and potential of spiking neural networks (SNNs) as a neuromorphic artificial intelligence architecture inspired by the behaviour of the human brain.

The project begins by presenting the state of the art in artificial intelligence, within a historical and technical framework that ranges from the earliest artificial neuron models to modern deep learning architectures. It describes the evolution of artificial neural networks, their main applications, and the limitations they currently present, especially regarding the computational and energy consumption of large-scale models. Based on this, SNNs are introduced as a possible evolution of ANNs, since they explicitly incorporate the temporal dimension and use event-based processing, which is closer to the biological functioning of neurons. This could give them an advantage in terms of energy efficiency and temporal processing capability.

Next, the technical and mathematical foundations of spiking neural networks are analysed. Existing neuron models are described, such as Hodgkin-Huxley, Integrate-and-Fire, and Leaky Integrate-and-Fire, as well as the role of synapses, temporal delays, and different information encoding strategies. The main training methods for SNNs are also examined, including local learning based on STDP, the use of surrogate gradients, and the conversion of a previously trained ANN into an SNN. This last strategy is the one used in the practical part of the project, as it makes it possible to take advantage of the maturity and relative ease of ANN training methods and transfer the learned weights to a spiking model.

The experimental part, presented next, is divided into two case studies. In the first one, a simple ANN is developed to classify points according to whether they are inside or outside a pseudo-randomly generated polygon. Once trained, this network is converted into an SNN using rate encoding, weight normalization, and temporal spike simulation. The results show that the ANN achieves an accuracy of 93.0\%, while the converted SNN reaches 90.1\%, with a relatively small loss in performance. These results confirm that ANN-to-SNN conversion is feasible and can reasonably preserve the behaviour of the original ANN model.

In the second case study, the same approach is applied to the MNIST dataset, a dataset consisting of images of handwritten digits. In this case, the starting point is a convolutional neural network (CNN) trained to classify images from 0 to 9, which is subsequently converted into an SNN. The CNN achieves an accuracy of 99.06\%, while the converted SNN reaches 93.19\%. Although the loss of accuracy is greater than in the geometric case, the results show that the conversion is also feasible in an image classification problem, a scenario that more closely resembles a real-world application. As a complement to visualise the results obtained, a web application has been developed that allows users to draw digits manually and compare the predictions of the CNN and the SNN.

Finally, the project analyses the possible economic, environmental, and social impact that the implementation of these types of networks could have. SNNs are not presented as an immediate replacement for conventional neural networks, but rather as a technological pathway with great potential in scenarios where energy consumption is critical. Their event-based nature can reduce computational activity and, therefore, electrical consumption and heat generation, which could help reduce the operating costs and environmental impact of artificial intelligence systems. The results obtained therefore indicate that SNNs are a promising alternative, especially if combined with specific refinement processes and neuromorphic hardware capable of exploiting their advantages.
\clearpage

\tableofcontents
{
	\let\cleardoublepage\clearpage
	\listoffigures
}

\section{Prefaci}
Aquest treball neix del meu interès per la intel·ligència artificial i, alhora, per l’electrònica i el disseny de hardware. Durant els darrers anys, la IA ha adquirit una importància creixent tant en l’àmbit tecnològic com en l’industrial, i s’ha convertit en una eina cada vegada més present en aplicacions de tota mena. Aquest progrés però, també ha posat de manifest alguns reptes importants, especialment relacionats amb l’elevat consum energètic dels models actuals, la necessitat de grans infraestructures de càlcul i l’impacte ambiental associat, incloent-hi no només l’electricitat consumida, sinó també l’aigua necessària per a la refrigeració d'aquests centres de dades.

Per aquest les xarxes neuronals de polsos, o Spiking Neural Networks (SNNs), m’han semblat una línia d’estudi especialment interessant perquè connecten la intel·ligència artificial i el hardware. A diferència de les xarxes neuronals artificials convencionals, les SNNs treballen mitjançant polsos discrets en el temps i modelitzen el comportament de les neurones d’una manera més propera a un circuit electrònic, amb variables com el potencial de membrana, el llindar de dispar, la fuita o la integració temporal del senyal. Això fa que el seu estudi no sigui només un problema de programació o d’aprenentatge automàtic, sinó també una aproximació a una arquitectura computacional amb una base física i electrònica.

Així doncs, aquest projecte m’ha permès apropar-me al món de la intel·ligència artificial des d’una perspectiva que encaixa amb el meu interès pel disseny de sistemes electrònics. L’objectiu d'aquest no és plantejar les SNNs com una substitució immediata de les xarxes neuronals tradicionals, sinó estudiar-ne els fonaments, entendre el seu funcionament i comprovar fins a quin punt poden reproduir el comportament d’una ANN. D’aquesta manera, el treball combina una part teòrica amb una part pràctica i serveix com a punt de contacte entre el software d’intel·ligència artificial i les possibles implementacions neuromòrfiques en hardware.

%----------------------------------------------------------------------------------------------
%----------------------------------------------------------------------------------------------
\section{Introducció}\label{introducció}
Durant els darrers anys, la inteligència artificial a patit un "boom" sense precedents, i s'ha convertit en una de les tecnologies més rellevants i prometedores del nostre temps, amb aplicacions en àmbits diversos, com la visió per computador, el processament del llenguatge natural, l'automatització industrial i la robòtica, o l'anàlisi de dades. Gran part d’aquest progrés ha estat possible gràcies al desenvolupament de xarxes neuronals artificials cada vegada més grans i complexes, capaces d’aprendre patrons a partir de grans volums d’informació. Tot i això, aquest creixement també ha comportat nous reptes, especialment pel que fa al consum energètic, la necessitat de hardware especialitzat i la dificultat d’implementar aquests models en dispositius amb recursos limitats. Les xarxes neuronals de polsos, o Spiking Neural Networks (SNNs), apareixen com una alternativa inspirada en el funcionament del cervell humà, basada en la transmissió d’informació mitjançant polsos discrets en el temps. Aquest treball se centra en l’estudi d’aquest tipus d’arquitectures neuromòrfiques, les quals podrien representar la solució als reptes anteriorment mencionats, amb l’objectiu d’entendre els fonaments, analitzar les diferències respecte a les xarxes neuronals artificials convencionals i comprovar experimentalment fins a quin punt una xarxa ANN entrenada pot ser convertida en una SNN capaç de conservar un comportament funcional similar.

\subsection{Objectius}
L’objectiu principal d’aquest treball és estudiar i desenvolupar una arquitectura d’intel·ligència artificial neuromòrfica basada en xarxes neuronals de polsos, analitzant tant els seus fonaments teòrics com la seva aplicació pràctica en tasques de classificació. Per aconseguir aquest objectiu, es plantegen els següents objectius específics:

\begin{itemize}
	\item Situar les SNNs dins del context històric de la intel·ligència artificial i de l’evolució de les xarxes neuronals.
	\item Estudiar les principals limitacions de les ANNs tradicionals, especialment en termes de consum energètic, escalabilitat i adequació a sistemes amb recursos limitats.
	\item Analitzar els fonaments tècnics de les SNNs, incloent-hi els models de neurona, la transmissió sinàptica, la codificació de la informació i els mètodes d’entrenament.
	\item Desenvolupar una ANN per a una tasca de classificació binària de punts dins o fora d’un polígon i convertir-la posteriorment en una SNN per verificar la viabilitat d'aquesta conversió.
	\item Comparar el rendiment de la ANN i de la SNN obtinguda, utilitzant mètriques com la precisió i la matriu de confusió.
	\item Aplicar el mateix concepte a un cas més proper a una aplicació real mitjançant el dataset MNIST, comparant una xarxa neuronal convolucional amb la seva versió convertida a SNN.
	\item Desenvolupar una demostració gràfica que permeti visualitzar i comparar el comportament dels models ANN i SNN de manera interactiva.
	\item Valorar el potencial de les SNNs des del punt de vista tecnològic, econòmic, ambiental i social.
\end{itemize}
\subsection{Abast del projecte}
L’abast d’aquest projecte inclou un estudi teòric de les xarxes neuronals de polsos i el desenvolupament pràctic de 2 models que permetin comparar el comportament amb el de xarxes neuronals artificials convencionals. En primer lloc, es presenta una revisió de la història de la intel·ligència artificial, des dels primers models de neurona artificial fins a les arquitectures modernes de deep learning.

Pel que fa al desenvolupament pràctic, el treball se centra en la conversió d’ANNs ja entrenades a SNNs. Aquesta conversió es duu a terme en dos escenaris diferents. El primer és un problema geomètric senzill, basat en la classificació de punts segons si es troben dins o fora d’un polígon. Aquest cas permet validar el procés de conversió en un entorn controlat i fàcilment interpretable. El segon escenari utilitza el dataset MNIST, format per imatges de dígits manuscrits, amb l’objectiu d’avaluar el comportament de la conversió en una tasca de classificació més propera a una aplicació real.

El projecte també inclou el desenvolupament d’una aplicació web de demostració, pensada per mostrar de manera visual les prediccions de la ANN i de la SNN sobre dígits dibuixats manualment. Aquesta eina té com a objectiu facilitar la interpretació dels resultats i permetre una comparació directa entre els dos tipus de models.

Queden fora de l’abast del treball la implementació física de la SNN en hardware neuromòrfic, el disseny d’un xip específic i l’entrenament directe d’una SNN des de zero mitjançant mètodes avançats com STDP o gradients substitutius. Tot i això, aquests aspectes es tracten a nivell conceptual per entendre el potencial futur de les SNNs.
%----------------------------------------------------------------------------------------------
%----------------------------------------------------------------------------------------------
\section{Context i marc teòric}

Per entendre com funciona la intel·ligència artificial actual, i especialment aquelles basades en estructures que simulen el funcionament del cervell humà com les xarxes neuronals de polsos (SNNs), és essencial mirar enrere i comprendre d'on provenen, i com han estat concebudes des dels seus inicis. La manera en com es va definir i enfocar la intel·ligència artificial en els seus primers anys, així com la posterior evolució dels models d'intel·ligència computacional i els diferents mètodes d'aprenentatge emprats al llarg de la història, continuen tenint un profund impacte sobre les tecnologies i els paradigmes utilitzats avui dia. La història de la intel·ligència artificial no és només una successió d'avenços tècnics, sinó també un reflex dels diferents models amb què la ciència ha intentat descriure el pensament, l'aprenentatge, la intel·ligència i la consciència humana. Explorar els orígens de la IA ens permet identificar els seus punts forts, els seus límits, i explica per què, després de dècades de desenvolupament en una direcció, estem tornant a idees fortament inspirades en la neurobiologia. Aquesta retrospectiva ens ajuda a situar les arquitectures neuromòrfiques dins d'un context evolutiu, i a entendre com aquestes poden donar una resposta a les limitacions dels enfocaments tradicionals utilitzats en l'actualitat.
%----------------------------------------------------------------------------------------------
\subsection{Orígens de la Intel·ligència Artificial}

Els primers esbossos del que avui dia entenem com a intel·ligència artificial sorgeixen l'any 1943 quan, amb la premissa de crear un "cervell electrònic", el neurofisiòleg Warren McCulloch i el matemàtic Walter Pitts publiquen \textit{A Logical Calculus of the Ideas Immanent in Nervous Activity} \cite{mcculloch1943}, un article que proposa un model matemàtic del sistema nerviós, a base d'una xarxa d'elements lògics simples, el que més tard es coneixerien com "Neurones Artificials". Aquestes  neurones, tot i fer-ho d'una manera simple basada en la lògica matemàtica, funcionen de manera similar a com avui dia funcionen les xarxes neuronals tradicionals, agafant "l'input" o entrada, realitzant-li una suma ponderada en cada neurona amb una funció llindar que filtra el moviment d'informació entre neurones, i generant "l'output" o sortida desitjada, tal com es mostra a la \autoref{fig:mcculloch_pitts}.
\begin{figure}[H]
	\centering
	\includegraphics[width=0.6\textwidth]{Neurona_de_McCulloch.png}
	\caption{Model de neurona artificial de McCulloch i Pitts.}
	\label{fig:mcculloch_pitts}
\end{figure}
Així doncs, la sortida que produeixen n d'entrades quedaria definida per la següent funció:
\begin{equation}
	\large f(\mathbf {x} )=h ( \sum_{n=1}^N (x_n w_n)  - \theta )
\end{equation}
Aquest article marca l'inici i estableix les bases de les xarxes neuronals artificials, obrint la porta a la idea que el pensament pot ser simulat per una màquina.

El primer en plantejar aquest paradigma és Alan Turing en el seu article \textit{Computing Machinery and Intelligence} \cite{turing1950}. En aquest, Turing planteja la pregunta de "Poden les màquines pensar?" introduint el que ara es coneix com el Test de Turing, un experiment de comportament que ell anomena \textit{Imitation Game} (joc d'imitació). Amb aquest experiment proposa que, si una màquina pot imitar les respostes d'un humà de manera suficientment convincent perquè un humà no pugui distingir de manera fiable d'una persona, aquesta màquina podrà ser considerada intel·ligent.
Turing també introdueix el concepte de \textit{self-teaching machines}, proposant la idea d'una màquina capaç de ser entrenada de manera similar a un infant humà, a base de recompenses i/o càstigs.

Basant en les idees fonamentals introduïdes al article de Turing, la universitat de Dartmouth marca el naixement formal de la intel·ligència artificial com a camp al 1956 amb el \textit{Dartmouth Summer Research Project on Artificial Intelligence} \cite{mccarthy1956dartmouth}, un programa proposat per John McCarthy, Marvin Minsky, Nathaniel Rochester i Claude Shannon un any abans, que es basa en la conjectura de que tots els aspectes de l'aprenentatge humà es poden descriure amb prou precisió perquè una màquina el simuli. Aquest projecte reuneix als principals pensadors de les matemàtiques, la informàtica, la neurociència i la ciència cognitiva de la època, i estableix per a dècades de recerca i innovació en aquest nou camp.

Basant-se en els conceptes fonamentals introduïts per Turing i formalitzats durant el Dartmouth Summer Research Project on Artificial Intelligence, el camp de la intel·ligència artificial entra en la seva primera fase important: la IA simbòlica, també coneguda com a \textit{Good Old-Fashioned AI} (GOFAI). Dominant la recerca en IA des de mitjans dels anys 50 fins a mitjans dels anys 90, aquest enfocament es centra en representar el coneixement mitjançant símbols llegibles pels humans i manipular-los mitjançant la lògica i la cerca heurística, tal i com s'il·lustra a la \autoref{fig:GOFAI}.
\begin{figure}[H]
	\centering
	\includegraphics[width=0.55\textwidth]{GOFAI.png}
	\caption{Il·lustració del funcionament de la Good Old-Fashioned AI (GOFAI).}
	\label{fig:GOFAI}
\end{figure}
Entre els primers programes importants basats en aquest enfocament s'inclou el \textit{Logic Theorist} o \textit{The Logic Theory Machine} \cite{newell1956logic} (1956), d'Allen Newell, Herbert A. Simons i Cliff Shaw, el primer programa en demostrar les teories proposades per Turing, aconseguint un sistema d'intel·ligència artificial capaç de raonament automatitzat demostrant teoremes matemàtics de \textit{Principia Mathematica} \cite{1927principiamatematica}.
Aprofitant l'èxit del programa, Newell, Simons i Shaw desenvolupen el \textit{General Problem Solver} (GPS) \cite{newell1959gps}, un sistema més versàtil que pretén resoldre una àmplia gamma de problemes mitjançant les estratègies heurístiques. Aquests esforços no només demostren que les màquines poden realitzar tasques tradicionalment associades amb el raonament humà, sinó que també inspiren futurs projectes de recerca i porten una creixent atenció per part del món acadèmic i del públic al camp emergent de la intel·ligència artificial. \newline

Cimentant-se en els fonaments establerts per les primeres instàncies d'IA simbòlica i el treball teòric sobre neurones artificials realitzat per McCulloch i Pitts, la dècada dels 50 veu també l'aparició de màquines d'autoaprenentatge, programes amb la capacitat d'adaptar-se basats en experiències en lloc de basar-se únicament en regles fixes com les màquines tradicionals. El programa de dames d'Arthur Samuel (1952-1955) \cite{samuelcheckers} marca un punt d'inflexió en el camp, demostrant l'habilitat de les màquines per aprendre i millorar sense ser explícitament programades per això, de manera similar al que proposava Turing d'ensenyar a una màquina com un infant. Més o menys al mateix temps, el model \textit{Perceptron} \cite{rosenblatt1958perceptron} va ampliar les idees de les neurones McCulloch-Pitts a xarxes neuronals entrenables, destacant una via paral·lela al raonament simbòlic. Altres innovacions, com l'experiment de traducció automàtica de Georgetown-IBM (1954) \cite{georgetown1954translation}, la primera demostració d'una màquina capaç de traduir frases senceres del rus a l'anglès, els primers programes de joc i el sistema \textit{STUDENT} \cite{bobrow1964student} per resoldre problemes d'àlgebra, van demostrar una creixent diversitat en els enfocaments de la IA, combinant mètodes estadístics, simbòlics i heurístics. Aquests desenvolupaments marquen el canvi del camp de la intel·ligència codificada a mà a sistemes capaços d'aprendre, adaptar-se i millorar, trets característics de la intel·ligència artificial moderna.

%----------------------------------------------------------------------------------------------
\subsection{Naixement xarxes neuronals}
Les xarxes neuronals són models computacionals inspirats pel funcionament del cervell humà. Consisteixen en grups de neurones artificials anomenades nodes, agrupades per capes. Aquestes, a través d'un \textit{input} o entrada i mitjançant sumes ponderades de node a node, aconsegueixen processar informació. El primer model del que avui dia anomenariem xarxa neuronal apareix amb el model del Perceptró (\textit{Perceptron} \cite{rosenblatt1958perceptron}) a mitjans dels anys 50. Concebut l'any 1957 per Frank Rosenblatt, el perceptró és un model de neurona artificial que utilitza un classificador binari conegut com a funció llindar (\textit{threshold function}), una funció que utilitza un nombre d'entrades (x1...xn), pesos (w1...wn) i un valor llindar o biax (b), junt amb la funció de Heaviside, per determinar l'activació de la neurona. El seu funcionament és similar al de la neurona de McCulloch i Pitts \cite{mcculloch1943}, amb la diferència fonamental que, en aquest darrer model, els pesos associats a les entrades són fixos i definits manualment, mentre que en el perceptró aquests pesos són paràmetres entrenables que s'ajusten automàticament durant el procés d'aprenentatge. \newline

\begin{equation}
	\large f(\mathbf {x} )=h ( \sum_{n=1}^N (x_n w_n)  + b )
\end{equation}

On \(\sum_{n=1}^N (x_n w_n) + b = z\). D'aquesta manera, la neurona s'activa \((f(x) = 1)\) si \(z \ge 0\) i queda inactiva \((f(x) = 0)\) si \(z < 0\), tal i com s'il·lustra a la \autoref{fig:perceptron}. Aquesta regla simple estableix si la informació s'ha de transmetre a la següent neurona o no, repetint aquest procés fins a arribar a la sortida (el valor de la darrera neurona) desitjada.
\begin{figure}[H]
	\centering
	\includegraphics[width=0.6\textwidth]{perceptron.png}
	\caption{Model de neurona del Perceptró.}
	\label{fig:perceptron}
\end{figure}
Aquest model de neurona possibilita la creació d'una xarxa interconectada que permet per primera vegada que una màquina aprengués de manera autònoma un criteri de classificació a partir d'exemples, en lloc de dependre exclusivament de regles manuals escrites per humans. Concretament, possibilita ajustar pesos i biaix a partir de dades d'entrenament mitjançant una regla local d'actualització, i l'adaptació i millora amb experiència sense necessitat de reprogramació explícita. A més, opera de manera eficient i en línia, actualitzant pesos per cada exemple, cosa que el fa útil per a sistemes adaptatius en temps real. Al mateix any de la seva concepció, l'any 1957, Frank Rosenblatt simulà un model del Perceptró en un IBM 704 al laboratori aeronàutic de Cornell (\textit{Calspan}), on es va provar la utilitat d'aquest model, el que serveix per cridar l'atenció del departament de la marina dels Estats Units (\textit{ONR}), qui més tard finançaria el projecte per crear una màquina dissenyada específicament per executar aquests models, el \textit{Mark I Perceptron}. Aquest model, però, no era pas perfecte. A causa de la seva arquitectura simple, el perceptró només pot separar les dades amb un hiperplà, és a dir, una línia (en 2D), un pla (en 3D), etc, com es mostra a la \autoref{fig:linear_separability}. Per tant, no podia resoldre problemes no linealment separables, cosa que va resultar problemàtica, al no poder resoldre la funció lògica bàsica XOR (exclusiva O). A més, la seva capacitat d'aprenentatge era limitada, doncs només podia aprendre patrons molt simples. Marvin Minsky i Seymour Papert van assenyalar totes aquestes debilitats al seu llibre Perceptrons (1969) \cite{minsky1969perceptrons}, fet que va provocar que molts investigadors i finançadors perdessin l'interès en el camp de la intel·ligència artificial. Això va donar lloc a un període conegut com l'hivern de la IA, caracteritzat per la baixa inversió i la poca investigació i interés en aquest àmbit, el qual va allargar fins a mitjans dels anys 80.

\begin{figure}[H]
	\centering
	\includegraphics[width=0.75\textwidth]{linear-separability.png}
	\caption{Problemes linealment separables i no linealment separables. Font: \cite{signatureVerification2009}.}
	\label{fig:linear_separability}
\end{figure}

Després de la desil·lusió causada per les limitacions del perceptró simple, alguns investigadors van continuar explorant la possibilitat de construir xarxes neuronals més potents. Ja a finals dels anys 50 i principis dels 60, es va començar a parlar de xarxes amb múltiples capes, però no es disposava d'un mètode eficient per entrenar-les, cosa que en limitava enormement l'aplicació pràctica. Tot i aquestes dificultats, la idea va quedar latent durant dècades fins que, l'any 1986, David E. Rumelhart, Geoffrey Hinton i Ronald J. Williams van publicar \textit{Learning representations by back-propagating errors} \cite{rumelhart1986learning}, un article que redescobria i formalitzava l'algorisme de retropropagació de l'error (\textit{backpropagation}) com a eina per entrenar xarxes neuronals multicapa. Aquest moment dona una base sòlida i funcional a un model que fins llavors havia estat només una idea prometedora: el perceptró multicapa (MLP).

El perceptró multicapa consisteix en una xarxa de neurones artificials organitzades en diverses capes consecutives. Aquest inclou una d'entrada, una o més capes ocultes, i una de sortida. Com està il·lustrat a la \autoref{fig:MLP}, cada neurona d'una capa està connectada amb totes les neurones de la capa següent mitjançant pesos que, de manera similar al Perceptró simple, regulen la importància de cada connexió. Quan s'introdueixen dades a la xarxa, aquestes passen per les capes de manera seqüencial. Cada neurona calcula una combinació lineal de les seves entrades i hi aplica una funció d'activació no lineal, normalment la funció de Sigmoid \cite{hahnloser2000digital} o la funció ReLU \cite{glorot2011deep}, que permet al sistema capturar patrons complexos.
\begin{figure}[H]
	\centering
	\includegraphics[width=0.75\textwidth]{MLP.png}
	\caption{Arquitectura d'un perceptró multicapa (MLP).}
	\label{fig:MLP}
\end{figure}
El gran repte, però, era saber com ajustar els pesos per tal que la xarxa produís resultats correctes. La retropropagació va permetre solucionar aquest problema. Aquesta funciona de manera que, una vegada la xarxa produeix una sortida, es compara amb la sortida esperada i es calcula un error. Aquest error es propaga cap enrere des de la capa de sortida fins a les primeres capes, utilitzant el càlcul del gradient de l'error respecte a cada pes. Això permet modificar els pesos mitjançant un algoritme d'optimització com el descens del gradient, de manera que l'error es pugui anar reduint a cada iteració. Aquest procés es repeteix milers o milions de vegades, amb moltes dades diferents, fins que la xarxa aprèn a realitzar la tasca desitjada amb gran precisió. El flux de dades i l'actualització de pesos es pot veure esmentat a la \autoref{fig:Backprog}.

\begin{figure}[H]
	\centering
	\includegraphics[width=1\textwidth]{Backprog.png}
	\caption{Flux de dades i actualització de pesos en la retropropagació.}
	\label{fig:Backprog}
\end{figure}

Aquests avenços suposen una veritable revolució, doncs per primera vegada, les xarxes neuronals eren capaces de resoldre problemes complexos, no lineals, i amb moltes variables, cosa que va reactivar l'interès pel camp de la IA i va posar les bases del Deep Learning que avui coneixem.

%----------------------------------------------------------------------------------------------
\subsection{Machine Learning i Deep Learning}

El \textit{Machine Learning} (o aprenentatge automàtic) és una branca de la intel·ligència artificial dedicada a l'estudi i al desenvolupament d'algorismes estadístics que aprenen a partir de dades. Aquestes algorismes permeten a les màquines generalitzar la informació obtinguda per resoldre tasques i prendre decisions en tasques per a les quals no han estat explícitament programades.
El terme Machine Learning l'introdueix per primer cop Arthur Samuel a la dècada dels 50 per descriure programes o màquines capaces d'aprendre i millorar el seu rendiment mitjançant l'experiència. La primera demostració de que això era possible fou el ja mencionat programa de dames \cite{samuelcheckers} del mateix Arthur Samuel. Durant les següents dècades, es van desenvolupar models i tècniques, com ara el perceptró i la retropropagació, que establirien els fonaments teòrics del Machine Learning, i popularitzarien el seu ús per l'entrenament de models computacionals durant la dècada dels 70. El Machine Learning utilitza principalment tres diferents paradigmes d'aprenentatge automàtic.

L'aprenentatge supervisat (SL), on l'algorisme aprèn a relacionar una entrada amb una sortida basant-se en conjunts d'exemples etiquetats, amb la finalitat de que aquest algorisme sigui capaç de predir etiquetes o valors per a noves dades de sortida en base a la entrada. Aquest model és el més estable, però no permet modelar decisions seqüencials.

L'aprenentatge no supervisat o aprenentatge autosupervisat (UL o SSL), que a diferencia del supervisat, utilitza dades sense etiquetes per extraure estructures estadístiques o patrons de les dades. Aquests algorismes no tenen la mateixa capacitat per realitzar tasques concretes com els entrenats per aprenentatge supervisat i poden requerir d'un aprenentatge supervisat reduit per poder realitzar-les, però al no haver d'etiquetar les dades i relacionar-les manualment amb la sortida, són bastant més fàcils d'entrenar, el que redueix el cost econòmic i el temps de fer-ho, especialment per a grans bases de dades no etiquetades.

Per últim, l'aprenentatge per reforç (RL), un mètode d'aprenentatge es basa en la psicologia conductista. En aquest, un agent o màquina rep informació de l'entorn (entrada), i emet accions (sortida). Un interprete, ja sigui humà o una funció de recompensa, retorna noves observacions (entrades), i una recompensa. L'objectiu de l'agent serà doncs aprendre un patró que maximitzi la recompensa acumulada, actualitzant de manera contínua els seus paràmetres segons l'experiència obtinguda. Aquest model, tot i ser ideal per resoldre problemes seqüencials i iteratius, és més costós, i pot presentar problemes d'inestabilitat i/o no convergència, els quals poden rendir el model inutilitzable per alguns usos.
Cada paradigma té fonaments teòrics i aplicacions pròpies i ha evolucionat amb contribucions històriques que van definir les tècniques i els reptes actuals.

Malgrat l'estancament provocat per les limitacions del perceptró i la crítica de Minsky i Papert als anys 60, la introducció i formalització de mètodes d'entrenament per a xarxes multicapa (notablement la retropropagació a mitjans dels anys 80) i, més tard, la disponibilitat massiva de dades i capacitat de càlcul que apareixen amb la popularització de l'internet impulsen l'avenç i l'expansió del camp, donant lloc al que avui anomenem Deep Learning.
El Deep Learning (DL) és una subàrea del Machine Learning que utilitza xarxes neuronals profundes, models multicapa amb moltes capes ocultes, normalment centenars o milers, per aprendre representacions jeràrquiques de les dades. La principal avantatge d'aquest tipus d'aprenentatge i de l'ús de xarxes neuronals profundes és la seva capacitat d'aprendre característiques directament a partir de dades brutes, com ara dades no etiquetades i no estructurades, el que les fa ideals per al seu entrenament amb algorismes d'aprenentatge no supervisat, aconseguint resultats similars als que es poden aconseguir utilitzant algorismes d'aprenentatge supervisat amb xarxes més simples, sense la necessitat d'estructurar i etiquetar les dades. Això es fa possible gràcies a les xarxes profundes desenvolupades als anys 80, com el perceptró multicapa, i a algorismes d'entrenament eficaços com la retropropagació combinada amb variants del descens de gradient (SGD, Adam, RMSprop) desenvolupats paral·lelament. Aquestes arquitectures permeten la creació d'aplicacions que modelaràn les següents dècades de desenvolupament del camp d'estudi de la intel·ligència artificial. D'aquestes tecnologies, destaquen les CNNs, o xarxes neuronals convolucionals, les RNNs, o xarxes recurrents i variants, i més endavant els transformers.

Les CNNs, o xarxes neuronals convolucionals, són xarxes dissenyades especialment per processar dades amb estructura espacial o temporal, com ara imatges, vídeos o àudio, de manera eficient. Les CNNs, proposades per primer cop a l'any 1980 per Kunihiko Fukushima, amb el seu "neocognitron" \cite{Fukushima1980}, model que era tas basava en el seu model previ d'elements de llindar analògics per models de visió \cite{Fukushima1969} presentat al 1969, i utilitzades per primer cop a l'any 1989 per Yann LeCun \cite{LeCun1989}, amb el seu model LeNet-5, dissenyat per a la classificació de caràcters manuscrits, el primer exemples d'una xarxa convolucional aplicada amb èxit. Fukushima, basant-se en les observacions del neurocientífic David Hubel i el neurofisiòleg Torsten Wiesel sobre com certes neurones al cervell responen en exposar un individu a certs estímuls visuals específics, va proposar la idea que podria existir un model jeràrquic per processar la informació visual. El model que proposa utilitza capes de convolució, capes que apliquen filtres a l'entrada per extraure característiques locals. Funciona transformant un tensor d'entrada (per exemple, una imatge amb amplada x alçada x canals (com ara, els colors dels píxels)) en un conjunt de mapes d'activació (feature maps) mitjançant operacions lineals seguides d'una no‑linealitat. Primer, aplica els filtres, que són petites matrius de pesos, per exemple de 3x3, a parts del tensor d'entrada, i calcula el producte. Aquest valor representa un punt del mapa d'activació. Aquest mateix filtre es mou per calcular un altre punt de la entrada, fins a tenir tots els productes de tots els píxels de la imatge calculats, i amb tots aquests es forma el mapa. Aquest procés es repeteix per a cada un dels filtres, donant com a resultat diverses capes per cada neurona, cada una asociada a un filtre, com está il·lustrat a la \autoref{fig:CNN}. Després, a aquestes capes se li aplica un biaix i una funció d'activació per introduir no linealitat a cada una. D'aquesta manera, podem aconseguir capes amb filtres més grans, que puguin detectar grans canvis de color en una imatge, permeten determinar contorns de grans objectes, i capes amb filtres més petits, que permeten diferenciar entre petits canvis a la imatge, cosa que, combinant-les, li dona l'habilitat a la xarxa d'identificació d'objectes en una imatge. És per això que aquests tipus de xarxes són les principalment utilitzades en el món de la visió per computador.

\begin{figure}[H]
	\centering
	\includegraphics[width=1\textwidth]{CNN.png}
	\caption{Estructura bàsica d'una xarxa neuronal convolucional (CNN).}
	\label{fig:CNN}
\end{figure}

Les RNNs, o xarxes neuronals recurrents, apareixen també a principis dels anys 1980, a mans de Rumelhart que, juntament amb David McClelland i Geoffrey Hinton, desenvolupa un algoritme de retropropagació a l'any 1986, presentat a l'article \textit{Learning Representations by Back-Propagating Errors} \cite{rumelhart1986learning}. A l'article, Rumelhart introdueix les xarxes neuronals recurrents com una extensió natural de les xarxes neuronals convencionals com l'MLP (Multi Layer Perceptrón), amb l'objectiu de poder tractar dades seqüencials o temporals, és a dir, informació on l'ordre dels elements és rellevant. A diferència de les xarxes tradicionals, on el flux d'informació va en una única direcció, d'una entrada cap a una sortida, les RNNs tenen connexions recurrents que permeten que la sortida d'una neurona en un període de temps pugui influir en el seu estat en períodes posteriors, es a dir, que la sortida de cada neurona no només depèn de la entrada, sinó que també depèn de les entrades ha tingut anteriorment, com es pot veure a la \autoref{fig:RNN}. D'aquesta manera, la xarxa manté una mena de "memòria" interna que li permet emmagatzemar informació del passat i utilitzar-la per processar noves entrades.
\begin{figure}[H]
	\centering
	\includegraphics[width=0.7\textwidth]{RNN.png}
	\caption{Estructura bàsica d'una xarxa neuronal recurrent (RNN).}
	\label{fig:RNN}
\end{figure}
Aquesta capacitat de recordar el context fa que les RNNs siguin ideals per realitzar tasques com el reconeixement de parla, la traducció automàtica o la predicció de sèries temporals, on les dependències temporals i de context són essencials. En comparació amb les xarxes neuronals estàtiques, les RNNs ofereixen una representació més dinàmica i flexible de les dades, podent modelar relacions al llarg termini dins d'una seqüència. Tot i això, malgrat el seu potencial, els primers models de RNN presentaven dificultats en l'entrenament, com el problema del gradient que s'esvaeix, on al entrenar una xarxa neuronal amb el mètode de retropropagació, a mesura que les dades es mouen d'una capa a una altra es van fer cada cop més petites, enlentint o directament parant el procés d'entrenament, o del gradient que explota, on les dades creixen de manera exponencial entre capes, produint el mateix efecte que el gradient que s'esvaeix, fets que limiten l'aplicació pràctica fins a la introducció de variants més estables com les LSTM i les GRU a finals dels anys 1990.

%----------------------------------------------------------------------------------------------
\subsection{Xarxes neuronals artificials modernes (ANNs)}
Les xarxes neuronals han recorregut un llarg camí des dels seus inicis, evolucionant en complexitat i eficiència amb l'avenç de la tecnologia i la investigació. Malgrat aquesta evolució, el fonament de les xarxes neuronals s'ha mantingut majorment inalterat. Les xarxes neuronals actuals es basen en els mateixos principis en els que es basaven les primeres xarxes neuronals, com ara el perceptró\cite{rosenblatt1958perceptron}, i més tard els MLPs \cite{rumelhart1986learning}, les CNNs \cite{LeCun1989} i les RNNs. Aquestes xarxes estan estructurades en tres tipus de capes. La capa d'entrada rep les dades inicials, on cada neurona representa una característica específica dels "inputs". Les capes ocultes, que poden ser múltiples, són les responsables de processar aquestes dades i extreure característiques complexes a través de connexions ponderades. Cada neurona d'una capa oculta combina les entrades que rep, aplica un pes i una funció d'activació abans d'enviar la informació a la següent capa. Finalment, la capa de sortida presenta el resultat del model, que pot ser una classificació, una predicció numèrica o qualsevol altra forma de resultat.

Els pesos assignats a cada connexió determinen la influència de l'entrada sobre la sortida. A més, les neurones poden disposar de biaixos, que afegeixen una constant per assegurar que les sortides puguin ser ajustades independentment de les entrades. Les funcions d'activació són fonamentals per a la creació de no linearitats dins del model, permetent que les xarxes aprenguin relacions complexes. Algunes de les funcions d'activació més populars avui dia inclouen la ReLU (Rectified Linear Unit), que retorna 0 per a entrades negatives i la pròpia entrada per a entrades positives, la sigmoide, que comprimeix les sortides entre 0 i 1, i la tanh, que les comprimeix entre -1 i 1.

El procés d'aprenentatge d'una ANN moderna es realitza en diversos passos. Primer, es fa una propagació cap endavant, on les dades es transmeten de la capa d'entrada fins a la sortida, passant per les capes ocultes i aplicant les funcions d'activació. Després, es calcula l'error comparant la sortida obtinguda amb la sortida esperada. Finalment, es realitza la retropropagació, un procés mitjançant el qual els pesos i biaixos són ajustats per reduir l'error. Utilitzant tècniques com el gradient descendent, la xarxa actualitza els pesos de manera que les futures prediccions siguin més precises.

Des dels anys 90, les xarxes neuronals han experimentat una evolució significativa no tant per la tecnologia en si, sinó pel seu escalat en termes de complexitat i capacitat. Inicialment, aquestes xarxes podien tenir centenars de milers de paràmetres (o neurones), limitades per la disponibilitat de dades i la potència de càlcul de la època. Amb l'accés a grans quantitats de dades que sorgeix amb l'aparició i popularització d'internet, i a maquinaria més potent, com les unitats de processament gràfic (GPUs), les xarxes neuronals han pogut creixer fins a arribar als milers de milions de paràmetres, permetent un modelatge més precís i detallat de patrons complexos, com ara la parla. Aquest increment en l'escala ha facilitat el desenvolupament d'intel·ligència artificial més avançada, capaç d'aprendre de conjunts de dades massius i de realitzar tasques que abans eren impensables, com ara el reconeixement d'imatges a gran escala, la traducció automàtica precisa i la generació de contingut de forma creativa.

Les arquitectures convencionals com les CNNs i les RNNs continuen utilitzant‑se per les tasques pels quals van ser dissenyades (CNNs per visió, RNNs per dades seqüencials), però al mateix temps han aparegut alternatives sovint més eficients o millor adaptades a problemes moderns, com els transformers per NLP (processament de llenguatge natural) i seqüències llargues, variants d'object‑detection com YOLOv8 \cite{Jocher2023} per detecció d'objectes en temps real en visió per computador, GANs (Generative Adversarial Networks) \cite{Goodfellow2014} per generació de contingut sintètic com imatges o videos i arquitectures híbrides que combinen components (per exemple: CNN+transformer), que en molts casos superen o complementen les CNN/RNN tradicionals, i que poc a poc van reemplaçant a aquestes xarxes. En els últims anys han estat sobretot els transformers els que s'han imposat com l'arquitectura més utilitzada i sovint la més efectiva en moltes tasques (especialment en NLP, però també en visió i multimodalitat), en part per la popularitat que intel·ligències artificials de processament de llenguatge basades en aquesta arquitectura han guanyat en aquest període, com ara el ChatGPT o Claude.

Els transformers apareixen fa relativament poc, més concretament a l'any 2017, amb la publicació de l'article \textit{Attention Is All You Need} de Vaswani et al. \cite{VaswaniNIPS}. En aquest article, s'introdueixen els transformers com una nova arquitectura de xarxa neuronal basada exclusivament en mecanismes d'atenció, eliminant així la necessitat de recurrència o convolució, que fins al moment caracteritzen les arquitectures anteriors com les RNNs o CNNs. \newline
Els transformers es basen en el mecanisme d'atenció, en concret en el mecanisme d'autoatenció, el qual permet a cada element d'una seqüència ponderar la rellevància de tots els altres elements de la mateixa seqüència per generar la seva pròpia representació.

Aquest procés es duu a terme mitjançant capes d'atenció múltiple (multi-head attention layers), que permeten que el model aprengui diferents tipus de relacions i dependències entre paraules o unitats de dades de manera simultània. A diferència de les RNNs, que processen les dades seqüencialment, els transformers poden processar tots els elements en paral·lel, fet que redueix significativament el temps d'entrenament i millora l'escalabilitat del model.

A més del mecanisme d'atenció, els transformers incorporen una sèrie de components addicionals com les "positional encodings" o codificacions posicionals, que afegeixen informació sobre la posició relativa dels elements dins la seqüència, i les "feed-forward networks" o xarxes de prealimentació, capes completament connectades que operen sobre cada posició de manera independent. Aquesta combinació de components fa que els transformers siguin extremadament flexibles i adaptables a diferents tipus de dades i tasques.

Els avantatges dels transformers són nombrosos. Ofereixen un rendiment superior en tasques de processament de llenguatge natural (NLP), permeten entrenar models amb una quantitat massiva de dades i paràmetres, i són altament paral·lelitzables, cosa que redueix els costos computacionals en comparació amb altres arquitectures seqüencials. Aquestes característiques han impulsat el desenvolupament dels grans models de llenguatge (LLMs), com la sèrie GPT (Generative Pre-trained Transformer), BERT (Bidirectional Encoder Representations from Transformers) o T5 (Text-To-Text Transfer Transformer), que han revolucionat el camp de la intel·ligència artificial moderna.

Avui dia, els transformers constitueixen la base de la majoria dels avenços en IA, no només en processament de llenguatge, sinó també en visió per computador, bioinformàtica, generació de codi i àrees multimodals que combinen text, imatge, àudio i vídeo. La seva capacitat per aprendre representacions contextuals complexes i generalitzar coneixement a tasques diverses els ha convertit en una de les fites més significatives en la història de la intel·ligència artificial i el deep learning, marcant l'inici d'una nova era dominada pels models fundacionals i els sistemes d'aprenentatge a gran escala.

Una part substancial del progrés recent en xarxes neuronals, incloent-hi l'aparició i consolidació dels transformers, no pot atribuir-se únicament a innovacions en l'arquitectura dels models. Aquesta evolució ha estat possible gràcies a un cúmul de factors, entre els quals destaca la millora significativa de les eines de programació, la disponibilitat d'infraestructures computacionals d'alt rendiment i l'accés a volums massius de dades digitalitzades. En aquest context, els llenguatges de programació actuals han tingut un paper fonamental. Python, en particular, s'ha posicionat com el llenguatge predominant en "machine learning" no només per la seva sintaxi accessible, sinó també per la maduresa d'un ampli conjunt de biblioteques especialitzades que faciliten totes les etapes del cicle de vida dels d'aquests models. Eines com TensorFlow, Keras, Pytorch i JAX proporcionen eines d'alt nivell que permeten definir i treballar amb arquitectures complexes d'una manera molt més senzilla i estructurada, optimitzar sèries numèriques de manera eficient i la integració nativa amb GPU i TPU, fet que permet entrenar models de gran escala que haurien estat inviables fa només uns anys. Paral·lelament, l'existència de paquets orientats a la gestió i preprocessament de dades com NumPy, Pandas o Dask, ha fet possible organitzar, transformar i manipular volums ingents d'informació d'una manera sistemàtica i reproduïble, cosa que ha facilitat notablement l'entrenament d'aquests models. Tot plegat ha generat un entorn on el desenvolupament, experimentació i validació de models complexos és molt més accessible, accelerant la innovació i contribuint directament a l'expansió de mètodes d'aprenentatge profund cada vegada més sofisticats.

Amb l'aparició de noves arquitectures neuronals i tecnologies d'aprenentatge automàtic, també han sorgit mètodes més especialitzats per classificar i avaluar els models segons les necessitats i la naturalesa de cada tasca. En el cas de la classificació, avui dia s'utilitzen mètriques com la precisió de la resposta obtinguda, el recall, que mesura la capacitat del model per identificar correctament els elements positius d'una classe, i la F1-score, que epresenta la mitjana harmònica entre precisió i recall, equilibrant els dos aspectes quan cap d'ells no pot ser ignorat, que permeten analitzar diversos aspectes del comportament del model, especialment en situacions amb classes desequilibrades. Per a tasques de regressió, en canvi, és habitual emprar mesures com l'error quadràtic mitjà (MSE) o l'error absolut mitjà (MAE), que quantifiquen la magnitud de les discrepàncies entre les prediccions i els valors reals. També s'utilitzen eines més específiques, com la matriu de confusió en problemes de múltiples classes, l'AUC-ROC (Area Under the Receiver Operating Characteristic Curve) \cite{fawcett2006roc} en classificació binària, el BLEU (Bilingual Evaluation Understudy) \cite{papineni2002bleu} en traducció automàtica o el mAP (mean Average Precision) \cite{everingham2010pascal} en detecció d'objectes, cadascuna pensada per capturar aspectes concrets del rendiment. L'elecció d'aquestes mètriques és essencial, ja que proporciona una visió precisa sobre la qualitat del model i orienta les millores necessàries per adequar-lo als objectius de l'aplicació concreta.

Un exemple clàssic que il·lustra de manera integrada els principis exposats fins ara és el conjunt de dades MNIST, un conjunt format per imatges de dígits manuscrits. Aquest recurs, àmpliament utilitzat en entorns educatius i experimentals, permet observar com diferents arquitectures, des de perceptrons multicapa fins a xarxes convolucionals, aborden una mateixa tasca de classificació. Per exemple, un model inicialment simple pot constar d'una capa d'entrada de 784 neurones (derivades dels 28x28 píxels que te cada imatge), una o diverses capes ocultes i una capa de sortida amb 10 neurones corresponents a les classes possibles (números del 0 al 9). Mitjançant el procés d'entrenament basat en propagació cap endavant i retropropagació del gradient, la xarxa aprèn a extreure patrons visuals rellevants a partir de milers d'imatges etiquetades. En aquest treball aquest mateix conjunt s'utilitzarà de manera il·lustrativa per demostrar com funcionen exactament les xarxes prèviament tractades, i per exemplificar les avantatges, desaventatges i diferències que poden tenir les xarxes de polsos respecte a les xarxes neuronals convencionals ja comentades.

%----------------------------------------------------------------------------------------------
\subsection{Limitacions de les ANNs}

Les xarxes neuronals artificials, malgrat els grans avenços i l'eficàcia demostrada en nombroses aplicacions, presenten limitacions que condicionen el seu desplegament i la seva adaptabilitat a certs entorns. En primer lloc, el cost computacional i energètic associat a l'entrenament de models grans és molt elevat. L'entrenament de xarxes amb milions o milers de milions de paràmetres requereix centres de càlcul amb GPUs/TPUs especialitzades, hardware que requereix d'un gran consum d'electricitat, i un temps significatiu d'entrenament, cosa que suposa costos econòmics i ambientals notables. A la \autoref{fig:Energy} es pot veure com ha evolucionat aquest consum a mesura que han evolucionat els models per a realitzar progresivament tasques de més complexitat. Això també es tradueix en una barrera d'entrada per a organitzacions petites i investigadors amb recursos limitats.

\begin{figure}[H]
	\centering
	\includegraphics[width=0.7\textwidth]{Ai_energy_consumption.png}
	\caption{Evolució del consum energètic en l'entrenament de models d'IA des de 2012 fins a 2020. Font: \cite{nature2019energy}.}
	\label{fig:Energy}
\end{figure}

Una altra limitació és la manca de temporalitat natural en moltes arquitectures estàndard. Tot i que les RNNs i variants com LSTM/GRU van ser concebudes per tractar dependències temporals, moltes arquitectures modernes, com per exemple els transformers, requereixen de mecanismes addicionals per captar de forma explícita la dinàmica temporal en dades contínues o en fluxos en temps real, fet que pot complicar el disseny de models per aplicacions com control en temps real o sèries temporals amb dependències llargues i no estacionàries.

Pel que fa a la correspondència amb el cervell humà, les ANNs no reprodueixen la major part dels principis biològics del processament neuronal. El seu funcionament és una aproximació matemàtica inspirada en el cervell, però diferent en arquitectura. Això implica que, malgrat algunes capacitats anàlogues (com l'aprenentatge a partir d'exemples), les xarxes artificials no presenten les mateixes propietats de robustesa, plasticitat i eficiència que caracteritzen el cervell humà.

L'escalabilitat amb eficiència és un repte addicional. L'augment de grandària dels models (per exemple, en LLMs) acostuma a portar millors resultats, però la millora no és lineal respecte al nombre de paràmetres. Sovint cal afegir ordres de magnitud més de paràmetres per obtenir guanys incrementals de precisió. En models molt grans com els LLMs actuals, per millorar notablement la qualitat cal incorporar molts més paràmetres, fet que fa que el cost d'entrenament i el consum energètic augmentin de manera quasi‑exponencial respecte a les millores obtingudes. Això empitjora la relació cost-benefici i complica la gestió pràctica, fent necessari revisar dissenys, tècniques d'eficiència (poda, quantització, knowledge distillation) i estratègies d'entrenament per contenir costos i consum, i limita l'escalament d'aquests models passat de models massius.

Finalment, hi ha limitacions rellevants per a "edge computing" o computació perifèrica, on el processament de dades i la presa de decisions es fan pròxims a la font de les dades per reduir la latència, i sistemes embarcats. La majoria de models d'alt rendiment són massa grans o exigents per a dispositius amb recursos limitats (CPU de baixa potència, memòria reduïda, restriccions d'energia), cosa que obliga a recórrer a tècniques de compressió, quantització, poda o a dissenys específics "light‑weight" o de baix pes o cost computacional, que sovint impliquen una pèrdua de precisió o capacitat. Aquestes restriccions condicionen l'adopció de les ANNs en aplicacions mòbils, IoT (Internet of Things) i en entorns amb requisits estrictes d'energia i latència.

%----------------------------------------------------------------------------------------------
\subsection{Xarxes neuronals de polsos (SNNs)}
Les SNNs, o Xarxes neuronals de polsos, apareixen com una evolució de les xarxes neuronals artificials tradicionals amb l'objectiu d'apropar encara més el funcionament dels models computacionals al comportament real del cervell humà. La seva motivació principal és neuromòrfica, és a dir, s'inspiren directament en la manera com les neurones biològiques processen i transmeten la informació, mitjançant impulsos elèctrics discrets, coneguts com a spikes o polsos. A diferència de les ANNs convencionals, que operen amb valors continus i activacions instantànies, les SNNs introdueixen explícitament la dimensió temporal com a part fonamental del càlcul neuronal.

En aquest tipus de xarxes, la informació no es codifica únicament en la magnitud del senyal, sinó també en el temps en què es produeixen els polsos. Aquesta codificació temporal permet representar la informació d'una manera molt més eficient i expressiva, ja que l'instant d'emissió d'un impuls pot contenir tanta o més informació que el seu valor. Les neurones acumulen el potencial d'entrada al llarg del temps i només emeten un pols quan aquest supera un cert llindar, imitant així el comportament de les neurones biològiques. D'aquesta manera, el processament esdevé inherentment dinàmic i dependent del temps, fet que fa que les SNNs siguin especialment adequades per tractar senyals temporals continus, com ara dades sensorials, fluxos d'esdeveniments o informació procedent de sensors neuromòrfics.

En comparació amb les ANNs tradicionals, les SNNs presenten diversos avantatges destacables. En primer lloc, permeten un consum energètic molt inferior, ja que només es produeix activitat quan hi ha polsos, fet que les fa especialment atractives per a sistemes encastats de baix consum, i per models de gran tamany, on la utilització de ANNs convencionals significa un gran cost energètic. En segon lloc, ofereixen una representació temporal més rica de la informació, cosa que facilita la captura de patrons temporals complexos sense necessitat de mecanismes addicionals. A més, la seva naturalesa basada en esdeveniments permet un processament més eficient i proper al funcionament del cervell, tant en termes de paral·lelisme com d'escalabilitat.

Tot i aquests avantatges, les SNNs també presenten desafiaments importants. L'entrenament d'aquestes xarxes és més complex que en les ANNs convencionals, ja que els mètodes tradicionals com la retropropagació no es poden aplicar directament a causa de la naturalesa discreta dels polsos. A més, la implementació pràctica de SNNs sovint requereix hardware especialitzat, com ara xips neuromòrfics, per aprofitar al màxim els seus avantatges energètics i de processament.
Malgrat aquests reptes, les SNNs representen una àrea de recerca molt activa i prometedora dins del camp de la intel·ligència artificial i el processament neuromòrfic, amb aplicacions potencials en robòtica, sistemes embarcats, i interfícies cervell-màquina, entre altres.
%----------------------------------------------------------------------------------------------
%----------------------------------------------------------------------------------------------
\section{Fonament tècnics de les SNNs}

Els primers models de neurona que avui dia s'utilitzen per la modelització de xarxes neuronals de polsos sorgeixen a l'any 1907 a mans de Louis Lapicque i, posteriorment, l'any 1952 a mans d'Alan Hodgkin i Andrew Huxley com a models matemàtics i electrofisiològics per a la simulació del funcionament d'una neurona humana, no amb la intenció de crear una intel·ligència artificial basada en aquestes. Aquests models neixen en el context de la neurociència amb l'objectiu d'entendre com les neurones biològiques processen i transmeten informació mitjançant senyals elèctrics. No és fins dècades més tard que aquests models són recuperats i reinterpretats com a fonament de les xarxes neuronals de polsos (SNNs).

Dins dels múltiples models proposats al llarg del segle XX, dos han esdevingut especialment rellevants per al desenvolupament de les SNNs, el model de Hodgkin-Huxley, d'alta fidelitat biològica, i el model Integrate-and-Fire, en particular la seva variant \textit{Leaky Integrate-and-Fire} (LIF), àmpliament utilitzada en models computacionals per la seva simplicitat i eficiència.

\subsection{Model de neurona de Hodgkin-Huxley}

El model de Hodgkin-Huxley, publicat l'any 1952 \cite{hodgkin1952}, és el primer model quantitatiu capaç de descriure amb precisió el comportament elèctric d'una neurona biològica. Aquest model va ser desenvolupat a partir d'experiments amb l'axó d'un calamar i descriu la generació i propagació del potencial d'acció mitjançant corrents iònics dependents del voltatge a través de la membrana neuronal.

A nivell elèctric, la membrana neuronal es modela com un circuit equivalent format per una capacitat (la membrana lipídica) en paral·lel amb diverses conductàncies variables que representen canals iònics de sodi (Na\textsuperscript{+}), potassi (K\textsuperscript{+}) i un corrent de fuga (leak). La dinàmica del potencial de membrana $V(t)$ ve donada per la següent equació diferencial:

\begin{equation}
	C_m \frac{dV}{dt} = I(t) - g_{Na} m^3 h (V - E_{Na}) - g_K n^4 (V - E_K) - g_L (V - E_L)
\end{equation}

on $C_m$ és la capacitat de la membrana, $I(t)$ el corrent injectat, $g_{Na}$, $g_K$ i $g_L$ són les conductàncies màximes dels canals de sodi, potassi i fuga respectivament, i $E_{Na}$, $E_K$ i $E_L$ els seus potencials d'equilibri. Les variables $m$, $h$ i $n$ són variables dinàmiques que representen l'estat d'obertura dels canals iònics i segueixen equacions diferencials pròpies dependents del voltatge. El circuit equivalent es mostra a la \autoref{fig:hodgkin_huxley}.

Aquest model reprodueix amb gran fidelitat fenòmens biològics com el llindar d'activació, el període refractari i la forma exacta del potencial d'acció. Tot i això, el seu elevat cost computacional el fa poc adequat per a xarxes neuronals grans, motiu pel qual rarament s'utilitza directament en SNNs a gran escala.

\begin{figure}[H]
	\centering
	\includegraphics[width=0.7\textwidth]{hodgkin_huxley_circuit.png}
	\caption{Circuit equivalent del model de Hodgkin-Huxley amb canals iònics dependents del voltatge. Font: \cite{hodgkingrafic1952}}
	\label{fig:hodgkin_huxley}
\end{figure}

\subsection{Model Integrate-and-Fire (IF) i Leaky Integrate-and-Fire (LIF)}

El model Integrate-and-Fire va ser introduït originalment per Louis Lapicque l'any 1907 al seu article \textit{Recherches quantitatives sur l'excitation électrique des nerfs traitée comme une polarisation} \cite{lapicque1907}, i posteriorment refinat per Stein en el seu article \textit{A theoretical analysis of neuronal variability} \cite{stein1965} en la seva variant més utilitzada, el model \textit{Leaky Integrate-and-Fire}. Aquest model de neurona és una simplificació extrema del comportament neuronal, però captura els mecanismes essencials necessaris per al càlcul en SNNs.

A nivell elèctric, la neurona Integrate-and-Fire (IF) es pot modelar com un sistema que integra el corrent d'entrada al llarg del temps i genera un pols discret quan el potencial acumulat supera un determinat llindar. El corrent d'entrada $I(t)$ pot provenir d'una font externa, en el cas d'una neurona d'entrada de la xarxa, o bé correspondre a la suma ponderada dels polsos generats per les neurones de la capa anterior. Addicionalment, el model pot incorporar un terme de biaix, representable com una font de corrent constant $I_b$, que permet desplaçar el punt de funcionament de la neurona.

Aquest corrent s'integra mitjançant un condensador, de manera que la tensió associada al potencial de membrana augmenta proporcionalment a la càrrega acumulada. El valor de tensió resultant es compara amb una referència mitjançant un amplificador operacional actuant com a comparador, que defineix el llindar de d'actuació $V_{ref}$. Quan la tensió supera aquesta referència, el comparador genera el pols de sortida i s'aplica un mecanisme de reinici del potencial. El esquema elèctric conceptual d'una neurona IF es mostra a la \autoref{fig:if_circuit}.

\begin{figure}[H]
	\centering
	\includegraphics[width=0.6\textwidth]{if_circuit.png}
	\caption{Esquema elèctric conceptual d'una neurona Integrate-and-Fire (IF).}
	\label{fig:if_circuit}
\end{figure}

Des d'un punt de vista matemàtic, la dinàmica del potencial de membrana $V(t)$ en el model IF ve descrita per:

\begin{equation}
	C \frac{dV}{dt} = I(t) + I_b
\end{equation}

on $C$ és la capacitat equivalent, $I(t)$ representa el corrent d'entrada i $I_b$ el terme de biaix, si s'utilitza. En absència d'entrada neta, el potencial es manté constant, ja que el model no incorpora cap mecanisme de descàrrega.

Quan el potencial de membrana assoleix el llindar $V_{ref}$, la neurona emet un pols i el potencial es reinicia a un valor $V_{reset}$:

\[
	\text{si } V(t) \ge V_{ref} \Rightarrow
	\begin{cases}
		\text{emet pols} \\
		V(t) \leftarrow V_{reset}
	\end{cases}
\]

A la \autoref{fig:if_waveform} es pot observar una representació temporal del potencial de membrana i dels polsos generats per una neurona IF en resposta a un corrent d'entrada.

\begin{figure}[H]
	\centering
	\includegraphics[width=0.85\textwidth]{if_waveform.png}
	\caption{Evolució temporal del potencial de membrana $V(t)$ en una neurona IF i generació de polsos en superar el llindar $V_{ref}$.}
	\label{fig:if_waveform}
\end{figure}

Aquest mecanisme introdueix una no-linealitat temporal i converteix la sortida de la neurona en una seqüència discreta d'esdeveniments temporals. No obstant això, el model IF presenta una limitació rellevant. Qualsevol corrent positiu sostingut, encara que sigui molt petit, acaba provocant un dispar, ja que la integració és indefinida en el temps.

Per corregir aquest comportament, Stein introdueix la neurona Leaky Integrate-and-Fire (LIF), que amplia el model IF afegint un mecanisme de pèrdua del potencial de membrana. A nivell elèctric, aquesta extensió es correspon amb la incorporació d'una resistència en paral·lel amb el condensador, que modela les pèrdues naturals de la membrana i força el potencial a relaxar-se cap a un valor de repòs $V_{rest}$. El esquema elèctric conceptual d'una neurona LIF es mostra a la \autoref{fig:lif_circuit}.

\begin{figure}[H]
	\centering
	\caption{Esquema elèctric conceptual d'una neurona Leaky Integrate-and-Fire (LIF).}
	\label{fig:lif_circuit}
\end{figure}

La dinàmica del potencial de membrana en el model LIF es descriu mitjançant:

\begin{equation}
	\tau_m \frac{dV}{dt} = -(V - V_{reset}) + R(I(t) + I_b)
\end{equation}

on $\tau_m = RC$ és la constant de temps de la membrana i $R$ la resistència equivalent. A diferència del model IF, en absència d'entrada el potencial decau exponencialment cap a $V_{reset}$, evitant la integració il·limitada.

De manera anàloga al cas IF, una representació temporal del potencial de membrana i dels polsos associats, com la mostrada a la \autoref{fig:lif_waveform}, resulta especialment útil per evidenciar l'efecte de la fuita i la relaxació del potencial entre entrades consecutives.

\begin{figure}[H]
	\centering
	\includegraphics[width=0.85\textwidth]{lif_waveform.png}
	\caption{Evolució temporal del potencial de membrana $V(t)$ en una neurona LIF, mostrant l'efecte de la fuita i la generació de polsos. Font: \cite{lee2020spikebased}}
	\label{fig:lif_waveform}
\end{figure}

Aquesta extensió aporta una millor estabilitat temporal, una major robustesa davant soroll i un comportament més proper al de les neurones biològiques, tot mantenint una elevada eficiència computacional. Per aquest motiu, el model LIF s'utilitza de manera majoritària en la simulació de xarxes neuronals espigades i en implementacions sobre hardware neuromòrfic.
\newpage
%----------------------------------------------------------------------------------------
\subsection{Sinàpsis i transmissió de polsos en una SNN}

Tot i que els models de neurona IF i LIF defineixen la dinàmica del potencial de membrana i el mecanisme de dispar, el comportament d'una xarxa neuronal de polsos ve determinat principalment pel model sinàptic. Les sinapsis són les encarregades de convertir els polsos presinàptics en corrents d'entrada que modulen el potencial de les neurones, actuant com l'element intermedi que connecta l'activitat discreta de la xarxa amb la dinàmica contínua del potencial, es a dir, l'encarregat de traduir la informació dels polsos a les neurones.

En una SNN, cada pols presinàptic no provoca un increment del potencial de membrana, sinó que genera un corrent transitori. Aquest corrent pot modelar-se mitjançant una funció de resposta temporal, habitualment anomenada kernel sinàptic. D'aquesta manera, el corrent sinàptic total que rep una neurona en un instant donat pot expressar-se com:

\begin{equation}
	I_{syn}(t) = \sum_{i,j} w_{ij} \sum_{k} \kappa(t - t_{j}^k)
\end{equation}

on $w_{ij}$ és el pes de la sinapsi que connecta la neurona presinàptica $j$ amb la neurona postinàptica $i$, $t_{pols}^k$ representa els instants de dispar dels polsos presinàptics, i $\kappa$ és el kernel sinàptic que descriu la forma temporal del corrent generat per un sol pols.
Aquest kernel pot adoptar diferents formes, sent les més comunes la funció exponencial, definida com:
\begin{equation}
	\kappa_{exp}(t) = \frac{1}{\tau_{s}} e^{-\frac{t}{\tau_{s}}} u(t)
\end{equation}
On $\tau_{s}$ és la constant de temps sinàptica i $u(t)$ és la funció d'esglaó unitari o funció de Heaviside, que garanteix la causalitat del model, i la funció alfa, definida per:
\begin{equation}
	\kappa_{alpha}(t) = \frac{t}{\tau_{s}} e^{-\frac{t}{\tau_{s}}} u(t)
\end{equation}


La funció exponencial modela un corrent que augmenta bruscament en el moment del pols i decau exponencialment després, mentre que la funció alfa presenta un augment més gradual seguit d'una decaiguda, oferint una representació més realista de la resposta sinàptica biològica. A la \autoref{fig:synaptic_kernels} es poden observar les formes típiques d'aquests kernels.
\vspace{2mm}

\begin{figure}[H]
	\centering
	\includegraphics[width=0.9\textwidth]{synaptic_kernels.png}
	\caption{Resposta temporal de la conductància sinàptica per a (a) una sinapsi exponencial i (b) una sinapsi alfa. Font: \cite{compneuroprinciples2011}}
	\label{fig:synaptic_kernels}
\end{figure}

Aquests tipus de kernel introdueixen una memòria temporal a la sinapsi, ja que l'efecte d'un pols persisteix durant un cert interval de temps abans de decaure. Això implica que el potencial de membrana d'una neurona reflecteix no només el nombre de polsos entrants, sinó també la seva distribució temporal.

En el context de neurones LIF, el corrent sinàptic s'integra conjuntament amb el mecanisme de fuita o leak de la membrana, donant lloc a una dinàmica que depén de la interacció entre la constant de temps de la sinapsi $\tau_{s}$ i la constant de temps de la membrana $\tau_{m}$. Aquesta doble escala temporal permet diferenciar patrons de polsos que arriben amb diferències temporals petites.

A més de la forma temporal del corrent, moltes formulacions inclouen explícitament un retard sinàptic associat a cada connexió. Aquest retard simula el temps necessari perquè un spike es propagui des de la neurona presinàptica fins a la postsinàptica, i es pot incloure modificant l'expressió del corrent sinàptic com:

\begin{equation}
	I_{syn}(t) = \sum_{i,j} w_{ij} \sum_{k} \kappa(t - t_{j}^k-d_{ij})
\end{equation}

on $d_{ij}$ és el retard sinàptic de la connexió $j \rightarrow i$. La introducció de retards permet que la xarxa sigui sensible no només a la presència de polsos, sinó també a la seva posició temporal, fent possible implementar mecanismes de detecció de patrons temporals, sincronització i càlcul basat en l'ordre d'arribada dels senyals.

En una formulació discreta en el temps, com per exemple en simulacions, el corrent sinàptic pot representar-se mitjançant una variable d'estat que es va actualitzant:

\begin{equation}
	I_{syn}(t+1) = \beta I_{syn}(t) + \sum_j w_{ij} s_j(t - d_{ij})
\end{equation}

on $\beta = e^{-\Delta t/\tau_{s}}$ i $s_j(t) \in \{0,1\}$ indica la presència d'un pols en el pas temporal. Aquesta expressió mostra com els polsos entrants contribueixen al corrent de manera acumulativa i amb una memòria temporal finita.

En conjunt, el model sinàptic constitueix l'enllaç entre la dinàmica individual de les neurones i el comportament col·lectiu de la xarxa. Mitjançant la combinació de pesos sinàptics, kernels i retards de propagació, les SNNs poden entendre i utilitzar informació distribuïda en el temps i explotar la dimensió temporal com a part del càlcul neuronal, diferenciant-se així de les xarxes neuronals artificials convencionals.

%----------------------------------------------------------------------------------------

\subsection{Codificació i entrenament en xarxes neuronals de polsos}

En les SNNs, la informació no es representa principalment mitjançant valors continus, sinó a través de seqüències de polsos discrets en el temps, anomenats spikes. Això implica que tant les entrades com les sortides del sistema han de ser codificades temporalment, i que la dinàmica interna de la xarxa depèn de la interacció entre aquests polsos i la memòria temporal de les neurones. La manera com es realitza aquesta codificació és un element fonamental en el disseny d'una SNN i condiciona tant el comportament dinàmic com l'estratègia d'entrenament.

Un dels primers aspectes clau en les SNNs és la codificació de la informació. La codificació defineix com un senyal, ja sigui continu o discret, es transforma en un tren de polsos que pot ser processat per la xarxa. Una de les estratègies més simples i utilitzades és la codificació per taxa (rate coding), en la qual la informació es representa mitjançant la freqüència de dispar d'una neurona al llarg d'un tram de temps. En aquest esquema, un valor elevat es correspon amb una taxa de dispar alta, mentre que un valor baix genera pocs o cap pols. Aquest tipus de codificació té una interpretació intuïtiva i és robust davant el soroll, ja que petites variacions en el temps dels polsos no afecten significativament la taxa mitjana. No obstant això, el rate coding sol requerir trams relativament llargs per obtenir una estimació fiable de la taxa, fet que pot incrementar la latència i el cost computacional.
En la pràctica, la codificació per taxa es sol implementar generant trens de polsos segons un procés estocàstic, habitualment de tipus Poisson, on la probabilitat d'emetre un pols en un interval de temps discret és proporcional al valor de l'entrada normalitzada. Aquest enfocament presenta l'avantatge de ser senzill d'implementar i relativament tolerant al soroll, però requereix finestres temporals prou llargues perquè la taxa d'actuació estimada sigui representativa del valor original. Malgrat aquesta limitació, el rate encoding resulta especialment adequat en escenaris de conversió de pesos ANN-SNN, ja que preserva la semàntica de les activacions mitjanes del model original.

Una altra estratègia de codificació són les codificacions temporals, que aprofiten de manera més directa el temps d'ocurrència dels polsos. En aquests esquemes, la informació no es codifica en el nombre total de polsos, sinó en els instants exactes en què aquests es produeixen o en les relacions temporals entre diferents neurones. Un exemple representatiu és la codificació \textit{time-to-first-spike}, on el valor d'un estímul es reflecteix en el retard fins al primer pols després d'un instant de referència. Altres variants inclouen la codificació per latència, que consisteix a representar la informació mitjançant el temps relatiu d'actuació respecte a un esdeveniment o estímul concret, la codificació per fase respecte a una oscil·lació global, on els polsos es posicionen dins d'un cicle oscil·latori i la fase en què es produeix l'actuació conté la informació, o la codificació basada en l'ordre d'actuació de diferents neurones, en la qual el patró informatiu ve determinat per la seqüència relativa de dispars, independentment del temps absolut. Aquestes estratègies permeten una representació molt més eficient i ràpida de la informació, sovint amb un nombre reduït de polsos, però són més sensibles al soroll temporal i requereixen una sincronització més acurada dels senyals.
Les codificacions temporals s'utilitzen principalment en escenaris on la latència i l'eficiència són factors crítics, o quan la naturalesa del senyal ja és inherentment temporal. Són especialment rellevants en el processament sensorial basat en esdeveniments, com ara visió neuromòrfica amb sensors dinàmics de visió (DVS), processament auditiu o sistemes que han de reaccionar en temps real amb un nombre mínim d'operacions. En aquests contextos, el fet de poder prendre decisions basades en els primers polsos permet reduir significativament el temps de resposta i el consum energètic. A més, aquestes codificacions exploten de manera més plena la dinàmica temporal de les neurones, apropant el funcionament del model a certs mecanismes observats en sistemes neuronals biològics. No obstant això, el seu ús acostuma a implicar un disseny més precís de la xarxa i dels mecanismes de sincronització, ja que petites variacions temporals poden afectar de manera significativa la representació de la informació. A la \autoref{fig:rate_temporal_encodings} es pot observar una comparació entre la codificació per taxa i diferents esquemes de codificació temporal, il·lustrant com cada mètode representa la mateixa informació mitjançant diferents patrons de polsos.

\begin{figure}[H]
	\centering
	\includegraphics[width=0.8\textwidth]{rate_temporal_encodings.png}
	\caption{Comparació entre codificació per taxa i codificacions temporals. Font: \cite{EncodingsZhao2017}}
	\label{fig:rate_temporal_encodings}
\end{figure}

Tot i que la codificació per taxa i les codificacions temporals són les més utilitzades avui dia, existeixen altres esquemes de codificació menys comuns però rellevants en determinats contextos. Un exemple és el \textit{burst coding}, en el qual la informació es representa mitjançant ràfegues de polsos emesos en intervals de temps molt curts. En aquest cas, no només és important la presència del dispar, sinó també el nombre de polsos dins d'una ràfega i la seva separació temporal. Aquest tipus de codificació s'ha utilitzat per modelar fenòmens biològics associats a l'atenció o a la detecció d'estímuls rellevants, i pot aportar una major robustesa davant el soroll en entorns on els estímuls són transitoris.

Un altre esquema alternatiu és la \textit{delta modulation}, on els polsos no codifiquen el valor absolut d'un senyal, sinó els canvis o increments respecte a un estat anterior. En aquest enfocament, un pols s'emet quan la variació del senyal supera un cert llindar, la qual cosa el fa especialment adequat per a senyals amb dinàmiques suaus o per a sistemes orientats a l'eficiència energètica. Aquest tipus de codificació és habitual en sistemes basats en esdeveniments, ja que permet reduir significativament el nombre de polsos generats quan el senyal és estable.

\subsection{Entrenament de xarxes neuronals de polsos}

L'aprenentatge de les SNNs presenta un repte addicional respecte al de les ANN convencionals, ja que la funció de dispar de la neurona és inherentment no deribable, i per tant no es pot emprar el mètode de la retropropagació. Una de les aproximacions més estudiades per abordar aquest problema és l'ús de regles d'aprenentatge locals inspirades en la neurobiologia, com la \textit{Spike-Timing-Dependent Plasticity} (STDP). En l'STDP, els pesos sinàptics s'ajusten en funció de la diferència temporal entre els polsos presinàptics i postsinàptics. Si un pols presinàptic precedeix un pols postsinàptic dins d'una període de temps curt, el pes sinàptic tendeix a augmentar, reforçant una relació causal. En canvi, si l'ordre temporal és invers, el pes es redueix. Aquest tipus d'aprenentatge és completament local, ja que depèn només de la informació disponible a la sinapsi, i permet capturar correlacions temporals sense necessitat d'un senyal d'error global.

Tot i els seus avantatges en termes de plausibilitat biològica i eficiència, l'STDP presenta limitacions importants quan s'aplica a tasques supervisades complexes. Per aquest motiu, s'han desenvolupat extensions com l'STDP modulada per recompensa, en la qual les actualitzacions sinàptiques locals es combinen amb un senyal global que reflecteix l'èxit o el fracàs del comportament de la xarxa. Aquest enfocament permet aplicar les SNNs a escenaris d'aprenentatge per reforç, on la recompensa pot estar retardada en el temps.

En paral·lel, una altra línia d'investigació molt activa es centra en l'adaptació dels mètodes de retropropagació del gradient a les SNNs mitjançant l'ús de gradients substitutius (surrogate gradients). En aquest enfocament, la funció de dispar, que és discontínua, es manté durant la simulació directa de la xarxa, però durant la fase de retropropagació se'n substitueix la derivada per una funció suau artificial. Aquesta aproximació permet aplicar una retropropagació a través del temps, o \textit{backpropagation-through-time}, sobre la dinàmica temporal de la xarxa i entrenar arquitectures profundes de SNN de manera supervisada. Malgrat que aquest mètode sacrifica part de la plausibilitat biològica, ofereix una gran flexibilitat i ha demostrat ser molt eficaç en tasques de classificació, processament de sèries temporals i aprenentatge seqüencial.

L'elecció entre codificació per taxa o codificació temporal, així com entre aprenentatge local basat en STDP o entrenament amb gradients substitutius, depèn fortament del tipus d'aplicació i dels requisits del sistema. Les estratègies basades en taxa solen ser més estables i senzilles, mentre que les codificacions temporals permeten aprofitar millor el potencial de càlcul temporal de les SNNs. De la mateixa manera, els mètodes d'aprenentatge local són especialment adequats per a sistemes distribuïts i implementacions neuromòrfiques, mentre que els mètodes basats en gradients ofereixen una major capacitat d'optimització global en entorns supervisats.

En conjunt, aquests mecanismes de codificació, dinàmica sinàptica i aprenentatge constitueixen els fonaments tècnics de les xarxes neuronals espigades. La combinació adequada d'aquests elements permet adaptar les SNNs a una àmplia varietat d'aplicacions, des del processament sensorial basat en esdeveniments fins a la classificació i el control en temps real, tot aprofitant la dimensió temporal com a part intrínseca del càlcul neuronal.

Una altra estratègia rellevant per a l'entrenament de xarxes neuronals de polsos consisteix a entrenar prèviament una xarxa neuronal artificial convencional i, posteriorment, transferir els pesos i biaixos a una SNN. En aquest enfocament, la ANN s'entrena utilitzant mètodes estàndard de retropropagació i funcions d'activació contínues, com ara la funció ReLU, amb les dades d'entrada i sortida transformades per ser compatibles amb la SNN (entrades i sortides $\in {[0,1]}$), mentre que la SNN hereda els paràmetres apresos. Aquesta conversió es fonamenta en l'equivalència funcional entre les activacions mitjanes d'una ANN i la taxa de dispar d'una neurona de polsos, de manera que els pesos sinàptics poden conservar, almenys de forma aproximada, el comportament del model original.

\begin {equation}
z = \sum_{j} w_j x_j + b
\qquad \Longleftrightarrow \qquad
V(t+1) = \alpha V(t) + \sum_{j} w_j s_j(t) + b
\end{equation}


El procés de conversió requereix habitualment una normalització dels pesos i una definició acurada dels llindars de dispar de les neurones de polsos, amb l'objectiu d'evitar saturacions i garantir que les taxes d'actuació es mantinguin dins d'un rang operatiu adequat. Els biaixos de la ANN es poden interpretar com a corrents constants o desplaçaments del potencial de membrana en la SNN, integrant-se així en la dinàmica temporal del model. Un cop realitzada la transferència, la SNN pot operar sense necessitat d'entrenament addicional, o bé refinar-se mitjançant tècniques d'ajustament (fine-tuning) basades en mètodes específics de SNN com els ja mencionats.

Aquest enfocament presenta diversos avantatges. En primer lloc, permet aprofitar la estabilitat dels mètodes d'entrenament de les ANN, evitant les dificultats associades a la no diferenciabilitat del dispar de les SNNs. A més, facilita l'ús de SNNs en tasques on les ANN ja han demostrat un rendiment elevat, proporcionant una via directa per beneficiar-se del potencial estalvi energètic i del processament temporal propi de les xarxes de polsos. No obstant això, aquesta estratègia també presenta limitacions. La conversió sol requerir finestres temporals prou llargues perquè les taxes de dispar convergeixin cap a les activacions equivalents, la qual cosa pot incrementar la latència. A més, no totes les dinàmiques temporals pròpies de les SNNs són explotades en aquest enfocament, ja que el comportament final queda condicionat pel model d'ANN original. Per aquest motiu, l'entrenament per conversió ANN-SNN s'utilitza principalment com una estratègia intermitja entre els mètodes clàssics de deep learning i els enfocaments específics de SNN, oferint un compromís entre rendiment, simplicitat i eficiència computacional.
Per posar a prova l'eficàcia d'aquest tipus de xarxa, i més concretament de la conversió ANN-SNN, al següent apartat es descriu el desenvolupament pràctic d'una ANN entrenada per la classificació d'un conjunt de dades, la seva posterior conversió a SNN i l'avaluació del rendiment obtingut.

%----------------------------------------------------------------------------------------------
%----------------------------------------------------------------------------------------------
\section{Desenvolupament pràctic: Classificació amb ANN i SNN}

Per tal de posar a prova els fonaments tècnics exposats en el capítol anterior i avaluar la viabilitat de les SNNs en escenaris reals, la part pràctica d'aquest treball se centra en l'estudi comparatiu entre xarxes neuronals artificials convencionals i xarxes neuronals de polsos aplicades a una tasca de classificació. L'objectiu és analitzar fins a quin punt les SNNs poden reproduir el comportament funcional d'un model ANN ja entrenada, així com identificar els avantatges i limitacions que presenta aquesta transició.

En primer lloc, s'ha entrenat una petita xarxa neuronal artificial per actuar com a classificador de dades, utilitzant mètodes d'entrenament estàndard basats en la retropropagació del gradient. Aquest model ANN serveix com a punt de referència. Un cop obtingut el model, els pesos i biaixos apresos s'han utilitzat com a base per a la construcció d'una xarxa neuronal de polsos equivalent.

A partir d'aquest model inicial, s'ha dut a terme la transformació dels paràmetres de la ANN cap a una SNN, reinterpretant els pesos sinàptics i els termes de biaix dins d'un marc temporal basat en polsos discrets. La SNN resultant s'ha avaluat sota les mateixes condicions que el model original, considerant criteris com la precisió i l'exactitud en la classificació de les dades, així com l'eficiència computacional.

Finalment, els resultats obtinguts amb ambdós models s'han comparat de manera sistemàtica, amb l'objectiu d'analitzar les diferències de rendiment i de funcionament entre els dos. Aquesta comparativa permet valorar en quines condicions les SNNs poden esdevenir una alternativa viable a les ANN convencionals, especialment en contextos on el processament temporal o l'eficiència computacional són factors rellevants. D'aquesta manera, la part pràctica no només valida els conceptes teòrics presentats, sinó que també ofereix una visió aplicada del potencial de les SNNs en la indústria moderna.
%----------------------------------------------------------------------------------------------
\subsection{Generació i transformació de les dades d'entrenament}

Per poder portar a terme l'entrenament de l'ANN, primer hem de generar i adaptar les dades a utilitzar. Per fer-ho, s'ha utilitzat el programa de MATLAB, doncs degut a les seves eines gràfiques integrades que simplifiquen la manipulació geomètrica, i a la seva funció nativa \textit{inpolygon}, que permet, amb una sola crida, determinar si un punt cau dins d'un polígon arbitrari, evitant així la necessitat d'implementar algoritmes de detecció de punts en polígon des de zero, el fan l'entorn ideal per a aquesta tasca. A més, MATLAB facilita la generació aleatòria de vectors i matrius i la interpolació de vèrtexs, funcions essencials per la generació de les dades necessaries.

Així doncs, s'ha creat un script que agafa com a entrada la quantitat de vèrtex que ha de tenir el polígon, i la quantitat de punts del perímetre a retornar, i retorna aquesta quantitat p de punts en una matriu $2\times p$ d'un polígon pseudo-aleatòriament generat. Aquest no és realment aleatori degut a les limitacions inherents dels ordinadors, que són màquines deterministes, i per tant qualsevol generador basat només en programari produeix pseudo-aleatorietat i mai serà completament aleatori. Per a l'objectiu de generar polígons però, aquesta pseudo‑aleatorietat és suficient.
Un cop creat l'escript, s'executa per generar el següent poligon, que es mostra a la \autoref{fig:polygon}.
\begin{figure}[H]
	\centering
	\includegraphics[width=0.5\textwidth]{PoligonClassificador.png}
	\caption{Polígon generat per l'script de MATLAB.}
	\label{fig:polygon}
\end{figure}

Una vegada generat el polígon que s'utilitzarà com a frontera de classificació, s'ha generat un altre script per generar les dades que s'utilitzaran per l'entrenament de l'ANN. Aquest script agafa com a entrada la quantitat de punts a generar i el polígon, expresat en punts del seu perímetre, i retorna una matriu 3xp, on cada columna representa un punt de coordenades (x,y) i la seva classe associada (0 o 1, depenent de si el punt cau dins o fora del polígon). La generació dels punts es realitza de manera pseudo-aleatòria, fent-se servir la funció nativa de MATLAB \textit{rand} per generar punts uniformement distribuïts dins d'un quadrat que engloba el polígon, i de la funció \textit{inpolygon} per determinar si cada punt cau dins o fora del polígon. Un cop obtinguda la matriu de dades, l'escript les guarda en un fitxer .csv que posteriorment serà carregat per l'entrenament de l'ANN, i genera una figura que mostra els punts generats, amb diferents colors segons la classe a la que pertanyen, com es pot veure a la \autoref{fig:training_data}.
\begin{figure}[H]
	\centering
	\includegraphics[width=0.6\textwidth]{ClassificacioDades.png}
	\caption{Dades d'entrenament generades, amb punts de classe 0 (fora del polígon) en vermell i punts de classe 1 (dins del polígon) en verd.}
	\label{fig:training_data}
\end{figure}

Amb el conjunt de dades exportat en format .csv, aquest es carrega a Python mitjançant el fitxer anomenat Dades.py, encarregat de dur a terme tot el preprocessament previ a l'entrenament de la xarxa neuronal. En primer lloc, aquest fitxer carrega les dades del .csv ajudant-se de la llibreria pandas, extreu les coordenades i les classes dels punts, i les ordena en columnes separades, de nom x i y per les coordenades, i Sortida per les classes. Tot seguit, les dades es converteixen al format numèric float32, format ampliament utilitzat en el deep learning degut a la seva eficiència computacional, i es normalitzen linealment les coordenades al rang [0,1], fent correspondre l'interval $x \in [-40,40]$ i l'interval $y \in [-20,20]$ a aquest mateix rang, intervals que es corresponen amb els utilitzats en la generació de les dades. En l'entrenament d'una ANN, i més concretament d'un MLP, com és el cas, no és necessari fer aquesta normalització, doncs els pesos i biaixos de la xarxa poden adaptar-se a les escales de les dades d'entrada. No obstant això, aquesta normalització es realitza per possibilitar la posterior conversió a SNN, ja que aquests tipus de xarxes requereixen que les entrades i sortides es trobin dins d'un rang específic, com ara $[0,1]$, ja que els polsos que generen les neurones de polsos es poden interpretar com a valors binaris. Un cop normalitzades, tant les dades d'entrada com les etiquetes es transformen en tensors de PyTorch, format que permet utilitzar les eines d'aquesta llibreria per a l'entrenament de la xarxa neuronal, i es divideixen de manera aleatòria en subconjunts d'entrenament, validació i test. Aquest pas evita que les mostres mantinguin l'ordre original del fitxer i garanteix una distribució més representativa dels exemples en cada subconjunt, alhora que permet avaluar de manera més fiable el rendiment del model en dades no vistes durant l'entrenament. Finalment, el fitxer Dades.py retorna els tensors corresponents als subconjunts d'entrenament, validació i test, que seran utilitzats en la fase d'entrenament de l'ANN.

%----------------------------------------------------------------------------------------------
\subsection{Creació, entrenament i avaluació de la ANN}

Un cop processades les dades, es procedeix a la definició i entrenament d'una xarxa neuronal artificial convencional, més concretament un MLP, que actuï com a classificador binari dels punts generats. La implementació s'ha dut a terme en Python utilitzant la llibreria PyTorch, escollida per la seva flexibilitat i facilitat d'ús en el desenvolupament de models de deep learning, així com per la seva àmplia comunitat i suport. Es comença definint una classe que hereta de \textit{torch.nn.Module}, que representa el model de xarxa neuronal. Aquesta classe inclou la definició de les capes lineals i les funcions d'activació, així com el mètode \textit{forward} que especifica com es processa la informació a través de la xarxa. Per aquest cas i tenint en compte que la finalitat del primer model és demostrar la viabilitat de la conversió a SNN, s'ha optat per una arquitectura senzilla, composta per una capa d'entrada de 2 neurones (corresponent a les coordenades x i y), una capa oculta de 6 neurones amb funció d'activació ReLU, i una capa de sortida de 1 neurona també amb funció d'activació ReLU. Per la tasca a realitzar per la ANN, una alternativa potser més adient hagués estat triar la funció d'activació sigmoide a la capa de sortida, ja que aquesta funció comprimeix les sortides al rang $[0,1]$, facilitant la interpretació de la sortida. No obstant això, s'ha optat per la funció ReLU per mantenir una coherència amb la posterior conversió a SNN, doncs les sortides de la ANN es poden interpretar com a taxes de dispar d'una SNN, i la modificació que aplica la funció d'activació Sigmoide sobre aquesta és no lineal, cosa que afectaria negativament al procés de conversió, mentre que la funció ReLU manté una relació més directa amb la conversió a SNN, només eliminant els numeros negatius, i deixant les sortides positives intactes. Un cop definida l'arquitectura del model, es procedeix a l'entrenament utilitzant un optimizador de tipus Adam i una funció de pèrdua de tipus Mean Squared Error (MSE). Per una tasca de classificació binaria, tornem a trobar que existeix una funció més adient, com és la Binary Cross-Entropy (BCE), però de nou obtem per una que facilita el procés de conversió, doncs a diferencia d'un model entrenat amb la funció BCE, un entrenat amb MSE no optimitza per una probabilitat logística “pura”, sinó per un valor escalar de resposta (en aquest cas 0 o 1), cosa que assimila més el seu comportament al d'un model SNN.
El procés d'entrenament es realitza durant un nombre determinat d'èpoques, durant les quals el model ajusta els seus pesos i biaixos per minimitzar la pèrdua en el conjunt d'entrenament. A aquesta quantitat arbitrària d'epoques se li aplica una funció d'aturada prematura, que mesura el rendiment del model en el subconjunt de validació després de cada època, el mostra, i si detecta que aquesta no millora, l'atura. Durant el procés d'entrenament, es registren en un arxiu separat anomenat $\texttt{entrenament\_log}$ les mètriques de rendiment, com ara la pèrdua i l'exactitud, tant en el conjunt d'entrenament com en el de validació, per a cada època. Es pot veure un exemple de les mètriques registrades en l'arxiu $\texttt{entrenament\_log}$ al \autoref{lst:logs_ann}.

\vspace{2mm}

\begin{lstlisting}[style=logs, caption={Fragment dels logs d'entrenament de la ANN}, captionpos=b, label={lst:logs_ann}]
Comencant entrenament de la xarxa.
Neurones d'entrada: 2
Neurones de capes ocultes: 6
Capes ocultes: 1
Neurones de sortida: 1
Taxa d'aprenentatge: 0.001
Decaiment de pesos: 1e-05
Funcio de perdua: Mean Squared Error Loss
Optimitzador: Adam
Epoques: 200

Epoca 01   train loss: 0.2143   val loss: 0.1987
Epoca 02   train loss: 0.1761   val loss: 0.1542
Epoca 03   train loss: 0.1328   val loss: 0.1194
...
Test Loss: 0.0412, Percentatge d'encert: 94.37 %
\end{lstlisting}

Un cop finalitzat l'entrenament, el model es guarda en un fitxer anomenat $\texttt{Pesos\_classificador.pth}$, que conté els pesos i biaixos apresos, i es procedeix a avaluar el rendiment del model en el subconjunt de test, obtenint una mètrica d'exactitud que serveix com a referència per a la posterior comparació amb la SNN convertida.

Per tal de comprovar de manera més intuïtiva i visual el comportament del model entrenat, s'ha desenvolupat un segon arxiu en Python anomenat $\texttt{Test\_model.py}$, destinat específicament a la fase de prova de la xarxa neuronal. Aquest script carrega el model prèviament entrenat i guardat al fitxer de pesos i restaurant l'estat dels paràmetres apresos durant l'entrenament. Un cop carregat, el model es posa en mode d'avaluació, cosa que garanteix que la xarxa treballi exclusivament en fase d'inferència i no modifiqui els seus pesos interns. A partir d'aquí, l'script permet a l'usuari introduir manualment les coordenades d'un punt concret del pla, i retrona la predicció que realitza l'ANN. Per fer-ho, les coordenades introduïdes es processen de la mateixa manera que durant la fase d'entrenament. Primer es normalitzen al rang $[0,1]$, mantenint la coherència amb el preprocessament aplicat al conjunt de dades original, i posteriorment es transformen en un tensor de PyTorch amb el format adequat perquè pugui ser processat pel model. A més, l'script incorpora el càlcul del valor màxim de sortida obtingut sobre el conjunt de validació, que s'utilitza com a factor de normalització per expressar la resposta del model d'una manera més interpretable, es a dir, en un rang de $[0,1]$, igual que amb les dades d'entrada. D'aquesta manera, en lloc de mostrar únicament la sortida bruta de la xarxa, es presenta una magnitud normalitzada que es pot interpretar com una mètrica aproximada de la confiança de que el punt es trobi dins del polígon. Aquesta normalització permet entendre millor la resposta del model, doncs es veu com, al apropar-se al límit del polígon, la sortida va disminuint progressivament, i com els punts que es troben clarament dins del polígon obtenen una resposta propera a 1, mentre que els punts clarament fora del polígon obtenen una resposta propera a 0, sent els punts de frontera els que obtenen respostes intermitges.
%----------------------------------------------------------------------------------------------
\subsection{Conversió a SNN}
Un cop entrenada i validada la xarxa neuronal artificial convencional, el següent pas consisteix en la seva conversió a una xarxa neuronal de polsos. Aquesta conversió s'ha implementat al fitxer $\texttt{SNN\_model.py}$, on es carrega el model ANN entrenat, se n'extreuen els pesos i biaixos, i es redefineix el càlcul de la xarxa dins d'un marc temporal discret basat en polsos. L'objectiu d'aquesta etapa no és tornar a entrenar el model des de zero, sinó conservar tant com sigui possible el comportament funcional de la ANN original i reinterpretar-lo. Per fer-ho, es parteix directament del classificador definit a $\texttt{Train\_model.py}$, amb arquitectura $2\rightarrow6\rightarrow1$, i se'n carreguen els pesos guardats després de l'entrenament.
La conversió es basa en l'equivalència de funcionament entre l'activació d'una neurona en una ANN i la taxa de dispar d'una neurona en una SNN. En lloc de treballar amb sortides contínues calculades en un únic pas, la SNN replica la mateixa transformació mitjançant una seqüència de passos temporals, en què les neurones acumulen corrent d'entrada, apliquen una dinàmica de memòria amb fuga i generen un pols quan el potencial de membrana supera un llindar determinat, com s'explica amb més detall a l'apartat anterior.
El primer element important de la conversió és la codificació de les entrades. En el model implementat no s'utilitza una codificació temporal complexa, sinó una codificació per taxa o \textit{rate encoding}. A cada pas temporal, cada coordenada d'entrada, ja normalitzada al rang [0,1], es transforma en un pols binari segons una regla estocàstica simple es genera un nombre aleatori uniforme i es compara amb el valor de l'entrada, de manera que com més gran és aquest valor, més probable és que la neurona d'entrada dispari en aquell instant. Això implica que una entrada propera a 1 generarà molts més polsos al llarg d'un interval de temps determinat que una entrada propera a 0. En conseqüència, la informació no es codifica en la magnitud directa del senyal, sinó en la freqüència relativa de polsos produïda durant els passos de simulació. En el codi, aquesta operació es realitza mitjançant l'expressió $\texttt{(torch.rand\_like(x\_input) < x\_input).float()}$, que constitueix el mecanisme de codificació d'entrada de la SNN. Al \autoref{lst:Codificacio_entrada} es pot veure un fragment del codi on es realitza aquesta codificació per taxa de les entrades, i com aquest modifica els paràmetres internals de la xarxa SNN.
\newpage
\begin{lstlisting}[style=codi, caption={Fragment de la codificació per taxa de les entrades, paràmetres interns i sortides}, captionpos = b, label={lst:Codificacio_entrada}, firstnumber=74]
	def forward(self, x_input):
              
        # --- CAPA 1 (Entrada -> Oculta 1) ---

        threshold= 1.0     # Llindar sortida
        decay = 0.9        # Leak (mantenim memoria)

        v1 = torch.zeros(self.w1.size(0))  # Inicialitzem potencials de membrana a zero
        v2 = torch.zeros(self.w2.size(0))  # Inicialitzem potencials de membrana a zero
    
        dispars_totals = 0
        
        for pas in range(self.t):
            
            # Input rate coding
            input_spikes = (torch.rand_like(x_input) < x_input).float()

            # CAPA OCULTA
            corrent1 = torch.matmul(input_spikes, self.w1.T) + self.b1
            v1 = v1 * decay + corrent1
            
            spikes1 = (v1 >= threshold).float()
            v1 = v1 - (spikes1 * threshold)  # Reset amb el threshold
            
            # CAPA SORTIDA 
            
            corrent2 = torch.matmul(spikes1, self.w2.T) + self.b2
            v2 = v2 * decay + corrent2
            
            spike_out = (v2 >= threshold).float()
            v2 = v2 - (spike_out * threshold)
            
            if spike_out > 0:
                dispars_totals += 1
\end{lstlisting}

El segon element important que apareix és la normalització dels pesos i dels biaixos de la ANN abans d'introduir-los a la SNN. En el codi, abans de crear la xarxa de polsos, es calcula l'activació màxima observada a la capa oculta i a la capa de sortida del model original sobre el conjunt de validació. Aquests valors, anomenats $\texttt{a1\_max}$ i $\texttt{a2\_max}$, s'utilitzen com a factors de normalització per dividir, respectivament, els pesos i els biaixos de la primera capa i els de la capa de sortida. Això permet reduir l'amplitud efectiva dels corrents injectats a les neurones de polsos i evitar que el potencial de membrana saturi massa ràpidament. En altres paraules, aquesta normalització serveix per adaptar l'escala d'activació de la ANN a l'escala de la SNN, on les neurones només poden emetre polsos quan superen un llindar fix. Si no es dugués a terme aquesta etapa, seria probable que algunes neurones disparessin de manera gairebé contínua o, a l'inrevés, que no arribessin mai a superar el llindar, degradant significativament el comportament del model convertit.
Un cop codificades les entrades i normalitzats els paràmetres, la dinàmica interna de la SNN es construeix amb dues capes de neurones de polsos. La primera rep els polsos d'entrada i calcula un corrent sinàptic a partir del producte entre els polsos i la matriu de pesos $W_1$, sumant-hi el vector de biaix $b_1$. El potencial de membrana de cada neurona es va actualitzant en el temps segons una dinàmica amb fuga, modelada per un factor \texttt{decay = 0,9}. Quan aquest potencial supera el llindar fixat a \texttt{threshold = 1,0}, la neurona emet un pols i el potencial es redueix restant aquest llindar (0,0 + biaix), implementant així un mecanisme de reset. La capa de sortida funciona de la mateixa manera, rebent com a entrada els polsos de la capa oculta i acumulant-los temporalment fins que genera els polsos finals de sortida. La resposta del model no és una probabilitat contínua com en el cas de la ANN, sinó la proporció de polsos emesos per la neurona de sortida al llarg de l'interval de temps, és a dir, $\texttt{dispars\_totals / t}$. En l'avaluació final, la implementació utilitza una finestra de \texttt{200} passos temporals. Al \autoref{lst:Normalitzacio_pesos} es pot veure un fragment del codi on es realitza aquesta normalització dels pesos i biaixos a les diferents capes de la SNN.

\vspace{3mm}
\begin{lstlisting}[style=codi, caption={Fragment de la implementacio de la classe SNN utilitzada i conversió de pesos i biaixos de l'ANN a l'SNN}, captionpos = b, label={lst:Normalitzacio_pesos}, firstnumber=57]
class SNN():
    def __init__(self, mlp_model, a1_max, a2_max, passos_temps=50):
        self.t = passos_temps
        
        # Extraccio de pesos i biaixos de l'MLP per a l'SNN:
        
        # Capa 1
        self.w1 = mlp_model.net[0].weight.data / a1_max  # Matriu 6x2
        self.b1 = mlp_model.net[0].bias.data / a1_max    # Vector 6
        
        # Capa Sortida
        self.w2 = mlp_model.net[2].weight.data / a2_max  # Matriu 1x6
        self.b2 = mlp_model.net[2].bias.data / a2_max    # Vector 1

\end{lstlisting}

Per avaluar la SNN convertida, el codi inclou una funció específica d'\textit{accuracy} o precisió, que evalua el conjunt de prova i, per a cada punt, calcula la sortida de la SNN, la compara amb un llindar de decisió i n'obté la classe predita. A diferència de la ANN, que en l'avaluació base utilitza un llindar de 0,5, per a la SNN un llindar de 0,59 per compensar la diferència en l'escala de sortida, doncs al haver utilitzat la funció ReLU a la ANN, l'output d'aquesta podia donar valors superiors a 1, i per tant el theshold de 0,5 no és exactament la meitat, i hem de compensar a la SNN per assimilar la seva precisió.
Aquesta mena de conversió acostuma a preservar raonablement bé el comportament de la xarxa en el nucli principal de cada regió de classificació, és a dir, en aquelles zones del pla on la ANN original produeix activacions clarament per sobre o clarament per sota del llindar. En aquestes regions, la separació entre classes és prou robusta perquè les fluctuacions introduïdes pel \textit{rate encoding} i per la discretització temporal tinguin un efecte limitat. Per contra, és esperable que les pèrdues de precisió apareguin sobretot a prop de la frontera de decisió, on petits canvis en la taxa de dispar, en el nombre de passos temporals simulats o en el llindar final poden alterar la classe assignada al punt. Aquesta degradació és coherent amb el fet que la SNN implementada utilitza una codificació estocàstica i una finestra temporal finita. La sortida d'aquesta és una aproximació mitjana al comportament de la ANN, però no una reproducció exacta punt a punt, especialment en les zones més ambigües del problema. No obstant això, en la pràctica, aquesta aproximació pot ser suficient per a moltes aplicacions, especialment aquelles on la robustesa i l'eficiència energètica de les SNNs compensen les petites pèrdues de precisió en les zones de frontera. Cal remarcar que, per utilitats on la precisió és crítica, és habitual aplicar després de la conversió un procés de reentrenament o ajustament (\textit{fine-tuning}) de la SNN, que permet refinar els pesos i biaixos dins del marc temporal de la SNN per recuperar part de la precisió perduda durant la conversió inicial. En aquest cas, però, s'ha optat per no realitzar aquest pas addicional, per tal de mantenir el focus en la comparativa directa entre el model ANN original i la SNN convertida sense entrenament addicional.
%----------------------------------------------------------------------------------------------
\subsection{Resultats i comparativa de rendiment}
Amb els dos models entrenats, s'avaluen ambdós sobre el conjunt de dades de validació amb l'objectiu d'obtenir una mesura del rendiment el menys esbiaixada possible. S'observa que la ANN obté un percentatge d'encert del 93,0\%, mentre que la SNN convertida assoleix un 90,1\%. A la \autoref{fig:confusion_matrix_ann} i a la \autoref{fig:confusion_matrix_snn} podem observar una matriu de confusió dels models de ANN i SNN respectivament que il·lustren els resultats d'aquesta evaluació, i mostren el tipus d'errors que cometen aquests.
\begin{figure}[H]
	\centering
	\includegraphics[width=0.52\textwidth]{Confusion_Matrix_MLP_Polygon.png}
	\caption{Matriu de confusió del model ANN per la classificació de punts.}
	\label{fig:confusion_matrix_ann}
\end{figure}
\begin{figure}[H]
	\centering
	\includegraphics[width=0.52\textwidth]{Confusion_Matrix_SNN_Polygon.png}
	\caption{Matriu de confusió del model SNN per la classificació de punts.}
	\label{fig:confusion_matrix_snn}
\end{figure}

Cal remarcar però, que la precisió de la SNN varia en funció del període de temps en que s'evalui, sent més precisa com més temps se li doni per evaluar l'entrada. Per al nostre cas, i com un ordinador no entén de temps, sinó de cicles de càlcul, hem fet servir 200 cicles. A més, la precisió de la SNN també depén d'un factor pseudo-aleatori, doncs la entrada que li donem no és una entrada real, sinó una transformació d'un punt a un encoding de polsos realitzats amb una funció aleatoria, el que implica que la sortida d'aquest model pot variar lleugerament entre execucions, afectant a la seva precisió. Tot i això la conversió ANN-SNN introdueix una pèrdua de rendiment relativament reduïda, d'aproximadament 2,9 punts percentuals, amb la principal font d’augment de l’error en la SNN sent una tendència més optimista en la classificació dels punts interiors, assignant com a interiors un nombre més elevat de punts que haurien de ser classificats com a exteriors, mentre que per el cas de la classificació dels punts realment interiors, els dos models es comporten identicament.

Tot i això, cal destacar que el comportament de la SNN per als punts realment interiors és pràcticament idèntic al de l’ANN. Per tant, la diferència principal entre ambdós models no es troba en la capacitat de reconèixer correctament els punts de dins, sinó en la major propensió de la SNN a incloure punts de fora dins la regió acceptada.. Aquesta diferència és coherent amb la pèrdua de precisió esperada al procés de conversió. Mentre que la ANN treballa directament amb valors continus, la SNN ha de representar aquesta informació mitjançant una codificació per taxa, una dinàmica temporal discreta i un recompte finit de polsos, fet que introdueix una certa aproximació respecte al comportament original del model. Per aquest motiu, és raonable que la xarxa convertida presenti una lleugera pèrdua de precisió respecte a la xarxa artificial entrenada directament.
Tot i així, des d'un punt de vista qualitatiu, la diferència obtinguda és prou petita per considerar que els dos models són raonablement similars en termes de precisió. Això indica que la conversió conserva de manera satisfactòria el comportament funcional de la ANN i que la SNN resultant continua essent capaç de classificar correctament la gran majoria dels punts.
Aquest resultat també apunta a una idea important. Si l'ANN de partida s'escalés a una arquitectura més gran amb l'objectiu d'aproximar encara més la seva precisió al 100\%, és esperable que la SNN convertida també millorés en la mateixa direcció. D'aquesta manera, seria possible obtenir una SNN amb una precisió molt elevada i, alhora, beneficiar-se de les avantatges pròpies d'aquest tipus de xarxes, com ara la seva naturalesa temporal i el seu potencial d'eficiència computacional, sense necessitat d'entrenar-la directament com una xarxa de polsos.
%----------------------------------------------------------------------------------------------
%----------------------------------------------------------------------------------------------
\section{Desenvolupament pràctic: Aplicació del dataset MNIST amb ANN i SNN com a simulació de cas d'ús real}

Per complementar l'estudi del rendiment dels diferents tipus de xarxa per una classificació geomètrica presentat a l'apartat anterior amb un exemple més proper a una aplicació real, en aquesta segona part pràctica s'ha obtat per utilitzar una xarxa neuronal més gran per a classificació d'imatges utilitzant el dataset MNIST. Aquest conjunt de dades està format per imatges en escala de grisos de dígits manuscrits del 0 al 9, amb una resolució de $28\times 28$ píxels, i es distribueix habitualment en 60.000 imatges d'entrenament i 10.000 imatges de prova. Es tracta, per tant, d'un problema de classificació senzill des del punt de vista visual, però prou representatiu per estudiar el comportament de xarxes neuronals.

\subsection{Justificació de l'ús de MNIST}
El dataset MNIST (Modified National Institute of Standards and Technology database) té el seu origen en els conjunts de dades de dígits manuscrits desenvolupats per NIST. En particular, es construeix principalment a partir de dades del NIST Special Database 19 \cite{NISTSD19}. Aquests conjunts eren un cúmul d'imatges de caràcters manuscrits pensats per a la recerca en reconeixement automàtic d'escriptura manuscrita, però mancaven de cap tipus d'estandarització, sent totes les dades de mides diferents, i sense distingir entre lletres i numeros. Yann LeCun i col·laboradors prepararen un nou conjunt de dades, l'MNIST, a partir d'una selecció, normalització i reformat de mostres procedents tant de NIST com del conjunt de dígits del USPS, un altre dataset de caràcters manuscrits no estandaritzat. Aquest procés tenia com a objectiu disposar d'un banc de proves estàndard, net i fàcilment utilitzable per a la recerca en reconeixement automàtic d'escriptura manuscrita. L'MNIST es va popularitzar especialment arran del seu ús en experiments de classificació de dígits com els associats a LeNet \cite{LeCun1989}, on servia per avaluar xarxes neuronals aplicades al reconeixement de documents i codis postals manuscrits.

Avui dia, l'MNIST continua sent un dels datasets més utilitzats en àmbits acadèmics i de recerca, especialment com a punt de partida en experiments de classificació d'imatges i en la validació inicial de noves arquitectures. La seva presència en eines àmpliament utilitzades com PyTorch o altres biblioteques d'aprenentatge automàtic mostren el seu paper com a conjunt de dades de referència dins de la docència i la investigació. Tot i que actualment existeixen conjunts molt més complexos i exigents, MNIST es manté com un estàndard útil per provar models de manera controlada, reproduïble i comparable amb treballs i tecnologies previes.

En aquest treball, MNIST s'ha fet servir per representar un cas d'ús més realista que el problema inicial de classificació de punts dins o fora d'un polígon. L'objectiu és treballar amb dades reals de naturalesa visual i amb una tasca de classificació coneguda, de manera que es pugui estudiar el comportament de la ANN i de la seva conversió a SNN en un escenari més proper a una aplicació real de reconeixement de patrons. D'aquesta manera, l'MNIST actua com un pont entre l'exemple geomètric senzill utilitzat com a demostració de viabilitat de la conversió en la primera part pràctica i un problema de classificació d'imatges més representatiu del tipus de tasques on habitualment s'apliquen xarxes neuronals.
%----------------------------------------------------------------------------------------------
\subsection{Preparació de dades i entrenament de la ANN}
Es comença la part pràctica carregant el dataset MNIST mitjançant les eines que ofereix la llibreria PyTorch, que inclou funcions específiques per descarregar i preparar aquest conjunt de dades. A diferència del cas d'estudi anterior, on haviem de normalitzar les dades manualment, en aquest cas, on les imatges són matrius de pixels en escales de grisos, on el valor de cada pixel ve representat com una escala de 0 a 1, sent el 0 el negre i el 1 el blanc, les dades ja es troben normalitzades de manera natural, i per tant no és necessari aplicar cap tipus de normalització addicional. Gràcies a la llibreria PyTorch, el dataset es carrega directament en format de tensors, i es divideix en subconjunts d'entrenament i de prova de forma automàtica, amb 60.000 imatges per a l'entrenament i 10.000 per a la prova, tal com es distribueix originalment.

Pel que fa a l'arquitectura del classificador, en aquest cas no s'ha utilitzat un perceptró multicapa simple (MLP) com en l'experiment anterior, sinó una xarxa neuronal convolucional. Aquesta elecció és fa degut a que les CNN estan especialment dissenyades per tractar dades amb estructura espacial, com és el cas de les imatges, i al tractar-se d'una extensió de la arquitectura MLP, on les dades es veuen processades a través de capes convolucionals abans d'entrar en la xarxa neuronal, la seva utilització no suposa cap desavantatge per a la conversió, provocant només una millora del rendiment de la ANN, i per tant una millor aproximació a l'hora de calcular els pesos de les neurones, els quals s'utilitzaran a la SNN. El model implementat, anomenat ClassificadorCNN, està format per dos blocs convolucionals seguits d'un classificador totalment connectat. El primer bloc aplica una convolució de 1 a 32 canals amb nucli 3x3, seguida d'una operació de max pooling que redueix la resolució espacial de 28x28 a 14x14, i d'una activació ReLU, escollida per els motius explicats a la primera part pràctica. El segon bloc aplica una segona convolució de 32 a 64 canals, novament seguida de max pooling, reduint la mida a 7x7, i d'una nova activació ReLU. Finalment, la informació s'aplana i es processa mitjançant una capa lineal de 3136 (7x7x64) a 128 neurones, una ReLU intermitja, i una capa de sortida de 128 a 10 neurones, corresponents a les deu classes possibles del dataset, és a dir, els dígits del 0 al 9. L'estructura d'aquest model es pot veure de manera més detallada al \autoref{lst:ClassificadorCNN}.

\vspace{3mm}
\begin{lstlisting}[style=codi, caption={Arquitectura del model ClassificadorCNN}, captionpos=b, label={lst:ClassificadorCNN}, firstnumber=12]
class ClassificadorCNN(nn.Module):
    def __init__(self, n_classes=10): # n_classes = 10 perque MNIST te 10 classes (digits del 0 al 9)
        super().__init__()

        # Bloc convolucional 1: extreu caracteristiques basiques (vores, textures)
        self.conv1 = nn.Sequential(
            nn.Conv2d(in_channels=1, out_channels=32, kernel_size=3, padding=1), # 1 canal d'entrada (escala de grisos) -> 32 filtres
            nn.MaxPool2d(kernel_size=2), # Reduccio espacial: 28x28 -> 14x14
            nn.ReLU()                  # Funcio d'activacio ReLU
        )

        # Bloc convolucional 2: extreu caracteristiques mes complexes
        self.conv2 = nn.Sequential(
            nn.Conv2d(in_channels=32, out_channels=64, kernel_size=3, padding=1), # 32 canals d'entrada -> 64 filtres
            nn.MaxPool2d(kernel_size=2), # Reduccio espacial: 14x14 -> 7x7
            nn.ReLU()                  # Funcio d'activacio ReLU

        )

        # Capes totalment connectades (classificador final)
        self.fc = nn.Sequential(
            nn.Flatten(),               # Aplanem el tensor: 64 x 7 x 7 = 3136 valors
            nn.Linear(64 * 7 * 7, 128), # Capa oculta totalment connectada
            nn.ReLU(),                  # Funcio d'activacio ReLU
            nn.Linear(128, n_classes)   # Capa de sortida: 128 -> 10 classes
        )
\end{lstlisting}

L'entrenament d'aquesta xarxa es realitza de manera similar a l'experiment anterior, utilitzant un optimizador Adam i paràmetres similars als utilitzats en l'experiment anterior doncs, tot i no estar entrenant la xarxa neuronal per la mateixa tasca, treballant un amb vectors i l'altre amb imatges, en ambdos experiments els models s'entrenen per tasques de classificació. El procés d'entrenament es duu a terme durant un nombre determinat d'èpoques, amb una funció d'aturada prematura basada en el rendiment del model en el subconjunt de validació, similar a l'apartat anterior, però utilitzant un nombre més reduït d'èpoques degut a la major quantitat de dades d'entrenament. Durant aquest procés, també es registren les mètriques de rendiment en un arxiu de logs, i un cop finalitzat, el model es guarda en un fitxer de pesos per a la posterior conversió a SNN.
%----------------------------------------------------------------------------------------------

\subsection{Conversió a SNN i configuració temporal}
Un cop entrenada la xarxa neuronal convolucional sobre el dataset MNIST, el següent pas consisteix en la seva conversió a una xarxa neuronal de polsos. Aquesta segueix, en essència, el mateix plantejament utilitzat en l'estudi anterior amb el classificador geomètric. En tots dos casos, es parteix d'una ANN ja entrenada i se'n transfereixen els pesos i els biaixos a una SNN que reinterpreta el càlcul de la xarxa dins d'un marc temporal discret basat en polsos. La codificació de les entrades es realitza mitjançant una codificació per taxa, on cada pixel de les imatges, ja normalitzat al rang [0,1], es transforma en un pols binari segons una regla estocàstica similar a la utilitzada en l'experiment anterior.

Per tal d'adaptar correctament la CNN al nou marc de càlcul, els pesos de cada capa es normalitzen abans de construir la SNN. Aquesta normalització es fa a partir dels pesos màxims d'activació observats a la CNN sobre el conjunt de validació. En aquest cas, però, per compensar la pèrdua d'intensitat del senyal que es produeix en transformar les entrades contínues en polsos binaris, no s'utilitza el màxim absolut, sinó que, a l'hora de calcular aquests pesos màxims, s'utilitza un percentil, en aquest cas del 96\%, de les activacions positives de cada capa. Aquest valor afecta doncs a l'escala dels pesos de la SNN, fent-los lleugerament més grans, i garantint així una senyal més forta i evitant que la conversió redueixi excessivament l'activitat interna de la xarxa. El mateix criteri s'aplica també als biaixos, que a més es divideixen pel nombre de passos temporals per repartir-ne la contribució al llarg de tota la simulació.

Pel que fa a la configuració temporal, la SNN funciona durant una finestra de 200 passos de temps, al llarg dels quals es van acumulant els polsos generats per l'entrada i per les diferents capes internes, alhora que decauen pel "leak" introduït. Aquest pas, de la mateixa manera que en l'estudi anterior, permet modelar el comportament temporal de la xarxa, permetent donar una resposta més ràpida baixant la quantitat de pasos a utilitzar, o una més lenta però fiable augmentant-los. La predicció final no es basa en logits continus, sinó en la taxa de dispar acumulada de les neurones de sortida, seleccionant finalment la classe amb més activitat. Igual que en l'estudi anterior, aquest procediment permet obtenir una aproximació funcional del model ANN original sense necessitat d'entrenar directament la SNN des de zero.
%----------------------------------------------------------------------------------------------

\subsection{Resultats i comparativa entre ANN i SNN}
Una vegada completada la conversió del model convolucional a una xarxa neuronal de polsos, s'ha procedit a comparar el rendiment dels dos models sobre el conjunt de prova de MNIST. Al evaluar-la en front del subconjunt de dades de prova, la CNN assoleix una exactitud del 99,06\%, mentre que la SNN convertida arriba al 93,19\%. Això suposa una pèrdua de precisió apreciable respecte al model original, però alhora mostra que la conversió conserva una part molt important de la capacitat de classificació de la xarxa artificial. Aquesta caiguda de rendiment és, en certa manera, esperable. A diferència de la CNN, que treballa directament amb activacions contínues, la SNN ha de representar la informació mitjançant polsos discrets distribuïts en el temps. Aquesta aproximació introdueix una simplificació del comportament del model original i, per tant, una certa pèrdua de fidelitat en la classificació final. Tot i això, els resultats obtinguts continuen sent prou bons per considerar que la xarxa de polsos manté un comportament funcional útil i que és capaç de classificar les imatges de manera relativament fiable.Per poder analitzar amb més detall el rendiment de cada model, es generen les matrius de confusió corresponents a la CNN i a la SNN, com a l'estudi anterior, il·lustrats a la \autoref{fig:confusion_matrix_cnn} i a la \autoref{fig:confusion_matrix_snn} respectivament.
\begin{figure}[H]
	\centering
	\includegraphics[width=0.65\textwidth]{Confusion_Matrix_CNN_MNIST.png}
	\caption{Matriu de confusió del model CNN per la classificació de dígits manuscrits.}
	\label{fig:confusion_matrix_cnn}
\end{figure}
\begin{figure}[H]
	\centering
	\includegraphics[width=0.65\textwidth]{Confusion_Matrix_SNN_MNIST.png}
	\caption{Matriu de confusió del model SNN per la classificació de dígits manuscrits.}
	\label{fig:confusion_matrix_snn}
\end{figure}
Aquestes matrius mostren que, mentre que la CNN cometeix molt pocs errors i tendeix a confondre classes específiques de manera molt limitada, la SNN presenta una major quantitat d'errors, especialment entre classes que tenen formes similars, com el 0 i el 9, el 4 i el 9, o el 8 i el 9, i quan la resposta temporal és dèbil o ambigua, la classe 0 tendeix a acumular més activitat relativa que la resta, donant com a resultat més errors on la predicció és 0 quan la classe real no ho era per a tots els digits. Aquesta tendència es correspón amb la pèrdua de precisió introduïda per la conversió, que afecta especialment a les zones de frontera entre classes, on les activacions de la CNN són més propenses a ser interpretades de manera errònia per la SNN a causa de la seva codificació per taxa i del recompte finit de polsos.

Cal remarcar però que, al igual que en el cas d'estudi anterior, es podria augmentar la precisió del model de polsos mitjançant un entrenament addicional posterior a la conversió. Aquest refinament permetria ajustar els pesos transferits des de la CNN per adaptar-los millor a la dinàmica temporal pròpia de la xarxa de polsos, sense necessitat d'entrenar la SNN des de zero. D'aquesta manera, la conversió ANN-SNN es manté com una estratègia útil per obtenir models de polsos amb un bon rendiment a partir de xarxes artificials ja entrenades, alhora que s'aprofiten els avantatges propis de les SNNs.
Cal recordar també que, per al cas de la SNN, la precisió obtinguda pot variar lleugerament entre execucions degut a la naturalesa estocàstica de la codificació per taxa, on les entrades es transformen en polsos binaris segons una regla aleatòria, i que a més, com s'ha comentat, la precisió de la SNN també depèn del nombre de passos temporals utilitzats per a l'avaluació, sent més precisa com més temps se li doni per evaluar l'entrada. En aquest cas, s'ha fet servir una finestra de 200 passos, però és possible que amb una configuració temporal diferent es pogués obtenir un rendiment lleugerament millor o pitjor.

A més de l'evaluació quantitativa, també s'ha realitzat una anàlisi qualitativa del comportament dels dos models, creant un script que retorna la predicció de la ANN i de la SNN per a una imatge concreta del conjunt de prova, així com el seu valor de sortida normalitzat, que es pot interpretar com una probabilitat aproximada de pertinença a la classe predita. Aquesta anàlisi permet observar com, en general, les imatges que la ANN classifica amb una alta confiança (sortida propera a 1) també obtenen una resposta forta per part de la SNN, mentre que les imatges més difícils o amb classificacions més dubtoses per part de la ANN (sortida propera a 0,5) tendeixen a obtenir respostes més intermitges per part de la SNN, reflectint així la pèrdua de precisió introduïda per la conversió i la naturalesa temporal i estocàstica del model de polsos.
%----------------------------------------------------------------------------------------------
%----------------------------------------------------------------------------------------------
\section{Demostració gràfica del comportament de la ANN i la SNN}
Amb l'objectiu de facilitar la prova pràctica dels models entrenats i disposar d'una eina interactiva de demostració, s'ha desenvolupat una petita aplicació web amb arquitectura client-servidor. Aquesta aplicació permet a l'usuari dibuixar una imatge de dígit manuscrit en una interfície gràfica des del navegador, i en enviar-la al servidor, aquest processa la imatge d'entrada i retorna les prediccions de la CNN i la SNN per a aquesta imatge concreta.

Per la part del servidor o "\textit{backend}", s'ha implementat en Python utilitzant Flask, un framework web, que és l'encarregat de carregar els models entrenats, processar les imatges d'entrada i retornar les prediccions de la ANN i la SNN quan ho demani el client.

El client o "\textit{frontend}", desenvolupat amb HTML i CSS, proporciona una interfície gràfica, més concretament un element de dibuix tipus canvas de $28 \times 28$ pixels, sobre el qual l'usuari pot escriure manualment un número. Un cop dibuixat el dígit, l'usuari pot premer el botó de classificació, que envia la imatge generada al servidor perquè sigui processada pels dos models, i retorna les prediccions dels dos models, permetent a l'usuari visualitzar les prediccions dels dos models de manera comparativa, fent-ne ús de barres de progressió per representar el percentatge de confiança d'encert. Un cop acabada la revisió dels resultats, aquesta permet tornar a analitzar un altre número, i així de manera successiva. Aquesta aplicació per tant, permet no només comprovar el rendiment dels models en casos concrets, sinó també observar com les diferències entre la ANN i la SNN es manifesten en la classificació de les imatges. A la \autoref{fig:app} es pot veure una captura de pantalla de l'aplicació web desenvolupada, on es mostra el canvas de dibuix, el botó de classificació i les barres de progressió que representen les prediccions de la ANN i la SNN per a una imatge concreta.

\begin{figure}[H]
	\centering
	\includegraphics[width=0.85\textwidth]{Web_interface.png}
	\caption{Captura de pantalla de l'aplicació web desenvolupada per a la demostració pràctica dels models ANN i SNN.}
	\label{fig:app}
\end{figure}
%----------------------------------------------------------------------------------------------
%----------------------------------------------------------------------------------------------

\section{Pressupost}

Des d'un punt de vista econòmic, el desenvolupament d'una SNN com la presentada en aquest treball és assumible com a prova de concepte, però és molt més costós quan es planteja una integració real en entorns industrials. Prenent com a referència el cost mig d'un enginyer a Espanya, uns 30 €/h com a referència més conservadora.

Per a un pilot industrial senzill basat en una ANN entrenada, la seva conversió a SNN, la validació, i la integració en un entorn de processament de dades com podria ser un \textit{datacenter}, es pot estimar una càrrega de treball d'entre 560 a 720 hores. Repartint-les entre anàlisi i especificació, desenvolupament del model ANN, conversió i ajust de la SNN, proves, monitoratge i desplegament, el pressupost de personal quedaria aproximadament entre 16.800 € i 21.600 € si s'utilitza la referència de 30 €/h. Amb una contingència del 15\% per incidències, documentació i iteracions addicionals, el cost puja a un interval aproximat de 19.320 € a 24.840 €. Aquestes xifres són coherents amb el desenvolupament de software avançat, però molt inferiors al cost d'entrenar o desplegar grans models des de zero.

Si, a més, es volgués portar aquesta idea a hardware dedicat, el cost augmentaria, però per a una xarxa petita com la del treball continuaria sent abordable a nivell de prototip. Una primera implementació en FPGA es podria fer amb plaques de desenvolupament comercials de baix o mitjà cost, amb preus que poden variar des dels 129 als 1.678 €. En un prototip petit, el cost del maquinari no seria el principal problema, sinó les hores de disseny HDL, verificació, síntesi, proves i integració. Per això, per a una implementació experimental d'aquest tipus, és raonable estimar un pressupost total d'entre 7.000 € i 10.000 € si es fa servir FPGA i una xarxa de mida limitada.

Des del punt de vista energètic, l'interès econòmic de les SNNs rau en el fet que el càlcul és escàs i basat en la pròpia estructura del circuit, el que li proporciona una notable millora en eficiència per tasques concretes. Intel afirma que el maquinari neuromòrfic Loihi 2 ha demostrat guanys d'eficiència d'ordre de magnitud en càrregues petites i de tipus edge, i IBM va reportar que el seu xip TrueNorth implementava 1 milió de neurones amb un consum d'uns 65 mW. Això mostra que, almenys en problemes adequats i en hardware específic, la idea neuromòrfica pot ser molt més eficient que l'execució convencional. Això dona peu a pensar que aquest tipus d'arquitectura podria ser interessant per entorns on el consum energètic és un factor crític, com ara dispositius mòbils, sistemes embeguts o aplicacions de computació a l'edge. En aquests casos, els guanys d'eficiència podrien justificar la inversió en desenvolupament i integració de SNNs, especialment si es tracta de tasques específiques on les SNNs poden excel·lir. Tanmateix, per a aplicacions de gran escala o de propòsit general, com els grans models de llenguatge, el balanç econòmic és més complex, doncs no s'ha implementat mai aquest tipus d'arquitectura per aquest pròposit. Tot i així podem fer una simulació dels possibles estalvis anuals en energia que aquest tipus de tecnologia podria aportar si s'apliqués per aquest propòsit.

A dia d'avui, els LLMs massius com ChatGPT, Gemini o Claude consumeixen quantitats enormes d'energia, i tot i que no hi ha dades públiques detallades sobre el seu consum energètic, es poden fer estimacions basades en les dades disponibles sobre el nombre de consultes i el consum per consulta. Per al cas de ChatGPT, OpenAI va indicar que al juliol de 2025 es processaven uns 18.000 milions de missatges per setmana \cite{openai_chatgpt_usage_2025}, i Sam Altman, el seu CEO, va afirmar que una consulta mitjana consumia aproximadament 0,34 Wh. Si es combina aquesta magnitud amb el preu mitjà final de l'electricitat al mercat espanyol el 2025, 83,45 €/MWh (és a dir, 0,08345 €/kWh), el consum anual quedaria entorn de 318 GWh/any, amb un cost elèctric aproximat de 26,6 milions d'euros anuals. En el cas de Gemini, Google ha publicat una estimació energètica per prompt de text en 0,24 Wh \cite{google_gemini_energy_2025}. També s'ha informat que l'aplicació Gemini havia superat els 750 milions d'usuaris mensuals. Com que Google no publica el nombre total de consultes/dia per a la seva app, una forma prudent de calcular-ne l'ordre de magnitud és assumir 100 prompts per usuari i mes. Amb aquesta hipòtesi, el consum anual seria d'uns 216 GWh/any, i el cost elèctric aproximat, aplicant el mateix preu de referència, d'uns 18,0 milions d'euros anuals. A partir d'aquí, es pot fer una estimació de l'estalvi potencial que podria aportar una arquitectura més eficient inspirada en SNNs. En un escenari conservador d'un 20\% d'estalvi energètic, ChatGPT podria reduir aproximadament 5,3 milions d'€/any, i Gemini 3,6 milions d'€/any. En un horitzó de 5 anys, això equivaldria aproximadament a 26,6 milions i 18,0 milions, respectivament. Si es considera un escenari més ambiciós del 50\% d'estalvi, els valors pujarien a 13,3 milions d'€/any per a ChatGPT, 9,0 milions d'€/any per a Gemini. Sumant els tres serveis sota aquestes hipòtesis, el cost energètic anual agregat seria de l'ordre de 45,2 milions d'euros, i un estalvi del 20\% representaria uns 8,9 milions d'euros anuals, mentre que un del 50\% s'enfilaria fins als 22,3 milions anuals. Cal remarcar però, que tot i que millores del 50\% en l'eficiència han demostrat ser possibles per a tasques determinades, no se sap si es podrien aconseguir en el cas de LLMs massius, i menys encara si es volgués mantenir la mateixa precisió i capacitat de resposta. Per això, aquestes xifres s'han de prendre com a estimacions molt aproximades, però que serveixen per posar en perspectiva el potencial econòmic d'aquesta arquitectura, i per justificar la inversió en el desenvolupament d'SNNs.

En conseqüència, la idea d'aplicar SNNs o hardware neuromòrfic no s'ha d'entendre tant com una solució immediata per substituir LLMs massius ja desplegats, sinó com una tecnologia amb potencial econòmic, sempre i quan els guanys d'eficiència observats en maquinari neuromòrfic puguin traslladar-se, ni que fos parcialment, a sistemes de gran escala.

\section{Impacte ambiental i social}

Seguint per la branca del apartat anterior, des del punt de vista ambiental, el principal atractiu de les SNNs és la seva potencial eficiència energètica. Una arquitectura de càlcul basada en esdeveniments pot reduir significativament l’energia necessària per fer funcionar models d'intel·ligència artificial, especialment si es compara amb grans models com les LLMs actuals. Si es mantenen les mateixes hipòtesis de consum utilitzades abans per a serveis de gran escala, només ChatGPT i Gemini sumarien aproximadament 534 GWh anuals de consum elèctric associat al funcionament d'aquests tipus de xarxes. Per fer una aproximació de l'impacte climàtic d'aquest consum, es pot utilitzar el factor d'emissió mitjà de CO$_2$ per kWh a Espanya publicat per la generalitat, que durant el 2023 \cite{gencat_mix_electric_co2}, últim any del que es tenen dades, va ser de 260g de CO$_2$ per kWh. Aplicant aquest factor al consum anual estimat d'aquests serveis, el seu impacte climàtic anual seria d'uns 138.840 tones de CO$_2$. En aquest context, un estalvi del 20\% en el consum energètic gràcies a l'aplicació de SNNs podria reduir les emissions associades en aproximadament 27.768 tones de CO$_2$ anuals, mentre que un estalvi del 50\% podria arribar a reduir-les en uns 69.420 tones anuals.

A aquesta reducció d’emissions se li suma un simple però important efecte físic: si els xips consumeixen menys electricitat, també dissipen menys calor. Això es deriva directament de la llei de Joule, de manera que una menor potència elèctrica es tradueix en una menor generació tèrmica. En entorns de datacenter, aquest aspecte és força important, doncs la refrigeració no és un element trivial. Segons l’IEA \cite{iea_energy_ai_2025}, els sistemes de refrigeració i control ambiental en centres de dades poden representar des d’un 7\% del consum total, fins a més del 30\% en instal·lacions menys eficients. Per tant, una reducció del consum del model no només estalvia energia de càlcul, sinó també una part de l’energia auxiliar necessària per evacuar la calor generada. Aquest menor escalfament, a més, també té implicacions ambientals indirectes relacionades amb l’aigua. Els centres de dades consumeixen aigua tant de manera directa, principalment per refrigeració, com de manera indirecta, a través de la generació elèctrica que els alimenta. El Lawrence Berkeley National Laboratory estima que els centres de dades dels Estats Units van consumir aproximadament 66.000 milions de litres d’aigua al 2023 \cite{shehabi2024datacenter}, i informes governamentals recents remarquen que el creixement de la IA està augmentant de forma significativa aquesta pressió sobre els recursos hídrics. Menys calor generada implica, en general, menys necessitat de refrigeració i, per tant, menys necessitat de l'ús d'aquest recurs. A més, en aquells casos en què la refrigeració es vincula a circuits d’aigua o a rebuig tèrmic cap al medi, com en el centres que utilitzen l'aigua de rius per la seva refrigeració, una menor dissipació de calor també ajuda a reduir l’impacte sobre l’entorn aquàtic. L’augment de la temperatura de l’aigua pot afectar negativament els ecosistemes fluvials i aquàtics. L’EPA assenyala que la temperatura de l’aigua afecta tots els organismes i que canvis tèrmics poden fer els hàbitats inhòspits per a determinades espècies \cite{epa_stream_temperature_2001}. De manera similar, estudis recents remarquen que l’aigua de rebuig més calenta pot incrementar riscos per a la biodiversitat aquàtica en reduir l’oxigenació del medi. Per tant, una infraestructura computacional més eficient no només redueix emissions i despesa energètica, sinó que també pot contribuir a minimitzar alguns dels efectes tèrmics i hídrics associats als grans centres de dades.

Des del punt de vista social, un efecte positiu d’aquest tipus de millora és la possibilitat de fer la IA més sostenible i més accessible. Si el cost energètic i operatiu d'aquests tipus de tecnologies baixa, el desplegament de sistemes intel·ligents pot ser més assumible no només per grans empreses, sinó també per entorns amb menys recursos, com centres de recerca petits o aplicacions per petits negocis. Això podria facilitar una adopció més equilibrada de la tecnologia i reduir parcialment la concentració de capacitat computacional en poques infraestructures gegants. Alhora, una menor demanda energètica també redueix la competència pels recursos elèctrics i hidràulics en territoris on la implantació de grans centres de dades ja genera tensió social i preocupacions ambientals.
%----------------------------------------------------------------------------------------------
%----------------------------------------------------------------------------------------------
\section{Conclusions}

Aquest treball ha tingut com a objectiu estudiar el funcionament de les xarxes neuronals de polsos (SNNs) i analitzar fins a quin punt poden esdevenir una alternativa viable a les xarxes neuronals artificials (ANNs) tradicionals. Per fer-ho, s'ha combinat una part teòrica, centrada en els fonaments històrics i tècnics de les xarxes neuronals en general, i més concretament de les SNNs, amb una part pràctica, on s'ha comprovat experimentalment la viabilitat de la conversió d'una ANN a una SNN en dos escenaris diferents: un problema de classificació geomètrica i un problema de classificació d'imatges basat en el dataset MNIST. Com a complement a aquests estudis, s'ha desenvolupat una aplicació web per a la comparativa visual i demostració pràctica del rendiment dels diferents models entrenats.

En primer lloc, el context i marc teòric ha permès situar les SNNs dins l’evolució de la intel·ligència artificial i de les xarxes neuronals. S’ha observat que, tot i que les ANNs modernes han assolit un nivell de rendiment molt elevat en una gran varietat de tasques, continuen presentant limitacions importants en termes de consum energètic, eficiència computacional i correspondència amb els mecanismes temporals del cervell biològic. Les SNNs apareixen doncs com una proposta d'evolució d'aquestes xarxes, ja que incorporen explícitament la dimensió temporal i treballen amb una representació de la informació basada en esdeveniments discrets. Aquesta característica les fa especialment atractives en entorns on l’eficiència energètica i el processament temporal són factors determinants.

Pel que fa a la part pràctica, el primer experiment, basat en la classificació de punts dins o fora d’un polígon, ha servit per estudiar el procés de conversió d’una ANN petita a una SNN. Els resultats mostren que la xarxa de polsos convertida és capaç de reproduir, amb força fidelitat, el comportament de la xarxa artificial original, amb una pèrdua de precisió relativament reduïda. Això confirma que la conversió directa ANN-SNN és funcional i conserva bona part del coneixement après pel model artificial, i legitimitza un estudi més profund d'aquest tipus de xarxa.

El segon experiment, basat en el dataset MNIST, ha permès traslladar aquest mateix plantejament a un problema més proper a una aplicació real. En aquest cas, el model de partida ja no és un perceptró multicapa simple, sinó una xarxa convolucional capaç de classificar imatges manuscrites. La seva conversió a SNN ha mostrat una caiguda de rendiment més visible que en el cas geomètric, però encara dins d’un marge prou bo per considerar la xarxa convertida com un classificador útil i relativament fiable. Aquest resultat demostra que la conversió ANN-SNN no queda limitada a problemes molt simples, sinó que també pot aplicar-se a tasques de reconeixement visual més representatives.

Una de les conclusions que es pot extreure d'aquest treball és que la qualitat de la SNN convertida depèn fortament de la qualitat del model ANN de partida. En tots dos experiments s’observa que la SNN manté una part molt significativa del comportament funcional de la xarxa artificial original, encara que amb una pèrdua inevitable associada a la discretització i a la codificació per taxa. Això suggereix que, si es parteix d’una ANN més gran, més refinada o més precisa, és raonable esperar que la SNN resultant també millori. Per tant, la conversió ANN-SNN es presenta com una estratègia robusta per obtenir xarxes de polsos competitives sense necessitat d’entrenar-les directament des de zero, aprofitant l’estabilitat i la maduresa dels mètodes d’entrenament per a xarxes més tradicionals.

També s’ha pogut comprovar que, tot i ser funcional, la conversió directa no és perfecta. La representació temporal basada en polsos introdueix una aproximació respecte al model continu original, i això es tradueix en una certa pèrdua de precisió. Això és especialment visible en els casos més complexos o més propers a les fronteres de decisió. Es per això que es pot concloure que el rendiment de les SNNs convertides podria millorar encara més mitjançant un procés posterior de refinament o entrenament addicional, ajustant els pesos transferits des de la ANN per adaptar-los millor a la dinàmica temporal del model de polsos.

Més enllà dels experiments realitzats i del seu rendiment extret per a tasques de classificació, el treball també apunta a que les SNNs tenen sentit no només com a model teòric inspirat en el cervell, sinó també com a línia tecnològica amb potencial pràctic real. El seu principal valor no és necessàriament superar avui dia les ANNs en precisió absoluta, sinó oferir una alternativa amb possibilitats d’eficiència energètica i d’implementació en hardware especialitzat. En un context on models molt grans, com els LLMs actuals, impliquen costos energètics i ambientals molt elevats, qualsevol arquitectura capaç de reduir significativament aquest consum esdevé d’interès.

Finalment, es pot concloure que els objectius del treball s’han assolit de manera satisfactòria. S’ha analitzat el marc conceptual de les SNNs, s’ha implementat un procés de conversió ANN-SNN en dos casos d’estudi, s’han comparat els resultats obtinguts i s’ha mostrat que la xarxa de polsos pot conservar una part molt elevada de la funcionalitat del model original. Tot això reforça la idea que les SNNs poden arribar a ser rellevants dins de l’evolució de la intel·ligència artificial, especialment si en el futur es combinen amb processos de refinament específics i amb hardware neuromòrfic capaç d’explotar els seus avantatges.
%----------------------------------------------------------------------------------------------
%----------------------------------------------------------------------------------------------
\section*{Agraïments}

Voldria expressar el meu agraïment, en primer lloc, al meu tutor, Álvaro Gómez Pau, per la seva orientació i disponibilitat durant la realització d’aquest treball. La seva ajuda no s’ha limitat únicament al desenvolupament del TFG, sinó que també m’ha servit de guia en altres aspectes tècnics relacionats amb el camp d'estudi i amb la manera d’afrontar problemes d’enginyeria. Aquest acompanyament ha estat molt valuós al llarg de tota aquesta etapa.

També voldria agrair a tots els professors del grau la formació rebuda durant aquests anys. Més enllà dels coneixements concrets impartits a cada assignatura, considero que han contribuït de manera important a la manera com entenc i abordo els problemes tècnics, i han estat una part fonamental del camí que m’ha portat fins a aquest treball.

Finalment, vull donar les gràcies a la meva família i al meu entorn proper pel suport, la paciència i la confiança que m’han transmès durant tot aquest procés. Sense aquest suport constant, arribar fins aquí hauria estat molt més difícil.
%----------------------------------------------------------------------------------------------
%----------------------------------------------------------------------------------------------

%INICI BIBLIOGRAFIA

\printbibliography[heading=bibintoc, title={Bibliografia}]

%----------------------------------------------------------------------------------------------


\end{document}